\chapquote{The Power of Love}{I never renounce my feelings. True Love can not be taken away by anyone. \\ \textemdash Sailor Moon }
\phantomsection
\label{cha:PowerOfLove}

\begin{multicols}{2}\raggedcolumns

\noindent A Princess is still human, and she is able to use all of humanity's varied and dangerous talents to her advantage. A Princess is also magical, and she is capable of focusing her hopes and beliefs into all sorts of potent advantages. 

In this chapter, we present all the unique abilities possessed by the Enlightened, from the magical Charms they wield to new Merits.


\section*{Echoes}
\phantomsection
\label{sec:Echoes}

The name you choose is a promise, and the Enlightened have chosen the name \textquotedblleft Noble." The Radiant know they are the example, they literally cannot help it, and so they have promised to be a good example. When a Noble is at her best people take notice, they are impressed and want to imitate her. When a Noble loses heart people notice and it scares them.  

This is magic, but only a little bit. A Princess' Inner Light makes her every action, good or bad, shine brighter. Human empathy does the rest. 

Whenever the Princess fulfills one of her Virtues roll Belief. On a success everyone who witnessed the Princess takes the Inspired condition which they may use on any roll to act in accordance with that Virtue. 

However if the Princess has a clear opportunity to fulfill her Virtue and chooses not to do so she must also roll Belief. If she fails the roll then everyone nearby takes the Shaken Condition unless they Succeed on a Composure + Supernatural Advantage roll. 

Generally, Echoes don't provide Princess's with any advantage over their antagonists or enemies. If a police officer forces a vigilante Champion to stand aside for the authorities he won't be Shaken by her hurt feelings. Her feelings will be more noticeable but the officer's own belief in the importance of due process is more than enough of a shield. This goes double for people who actually want the Princess to suffer, while some Princesses have witnessed their enemies having a \textquotedblleft my god, what am I doing" moment any Princess who expects or relies on her Echo protecting her from her enemies is an idiot. If you prefer more defined rules for when a Princess' Echo can help her against her antagonists: A character who triggers a Princess's Echo by acting in accordance with their Virtue or Aspirations (or equivalents) will not be Shaken by by it. 

Princesses have their Echo in their mundane identity. 

\section*{Transformation}
\phantomsection
\label{sec:Transformation}
The most essential expression of a Princess's power is the ability to Transform. She calls on the magic she's gathered to herself and assumes the form of her \textquotedblleft perfected self", an idealized version of herself unfettered by this world of sorrows. Although she can only maintain this immense power for a limited time, a Transformed Princess is a force to be reckoned with.

The first time a Princess Transforms is called her Blossoming. This is an unconscious act, usually in response to some great need or danger. After her Blossoming, a princess' Phylactery will manifest. This is the physical representation of the power sealed away inside of her, and may take the form of any small simple object. She must have her Phylactery with her later to Transform. When a Princess wishes to Transform, she has two options: she may either spend a Wisp or make a Transformation roll. If she spends a Wisp, she Transforms automatically as a reflexive action. The Transformation roll is Belief + Inner Light - Shadows. 

\begin{myindentpar}{10pt}
    \emph{Dramatic Failure:} The Princess's Phylactery discorporates. She must spend a Wisp and a Willpower point to resummon it, which takes a full round and requires her absolute concentration. 
    
    \emph{Failure:} The Princess fails to Transform and spends the full round. 
    
    \emph{Success:} The Princess Transforms as an instant action. 
    
    \emph{Exceptional Success:} The Princess Transforms as a reflexive action.
\end{myindentpar}

A Princess can return to her ordinary self at any time as a reflexive action. She must de-transform if she spends her last Wisp, or when she falls asleep. If a Princess tries to Transform while her wisp pool is empty she fails. Also, transformation grows harder to sustain as time passes. When a Princess reaches the end of a scene transformed, or after half an hour spent transformed, she rolls Belief + Inner Light: 

\begin{myindentpar}{10pt}
    \emph{Dramatic Failure:} The Princess returns to her mundane self and her Phylactery discorporates. She must spend a Wisp and a Willpower point to resummon it, which takes a full round and requires her absolute concentration. 
    
    \emph{Failure:} The Princess must spend a Wisp to remain transformed in the next scene, or the next half-hour, and future rolls to remain transformed take a -2 cumulative penalty until she ends a scene in her mundane self. 
    
    \emph{Success:} The Princess remains transformed at no cost. 
    
    \emph{Exceptional Success:} The Princess remains transformed, and need not make this check at the end of the next scene, or for a full hour. 
\end{myindentpar}

While transformed, a Princess has full access to all of her Charms and abilities. 

\subsection*{Phylacteries}

Every Princess has a Phylactery, a physical item that symbolizes exactly what being a Princess means to her. A Phylactery is also the source of a Princesses power to Transform, she must activate it, open it, apply it to herself, or at least touch it if she wishes to Transform, this limit affects anything described as a Transformation action. Many Phylacteries also Transform with the Princess, becoming larger and more ornate or becoming the Regalia linked to a Charm strongly associated with her duties (a Phylactery is simply a specialized piece of Regalia and share most of their rules). The rest remain as they are, but are incorporated into her Heraldry. 

A Phylactery has a Durability equal to Inner Light and a Structure equal to Belief, typically it will have a Size of 0 - 2. If her Phylactery is destroyed or lost, the Princess can recreate or summon it. She must spend a Wisp, a point of Willpower and concentrate for a full round. 

If another was to steal a Princess' Phylactery they would be unable to use it, but they would be able to target the Princess through an unbreakable Intimacy link with a strength equivalent to love. Fortunately any technique that would cut the connection Princesses use to summon her Phylactery would, by definition, destroy it.

\section*{Dual Identity}
\phantomsection
\label{sec:DualIdentity}

One aspect of the Hopeful's transformation is subtle but amazingly convenient: the mundane and Transformed selves always appear to be different people. Even if the two identities look as similar as Clark Kent and Superman differences in a Princess' posture and bearing, and powerful supernatural protection, will make the very idea that they are the same person laughable. A simple perception roll never reveals a Noble's two identities are connected, nor can a simple subterfuge or empathy roll uncover the Princess' deception when she claims to have only one identity. 

Learning a Princess' secret identity requires iron clad proof, such as witnessing the Princess Transform with your own eyes. If the proof is solid but not iron clad, such as finding the Princess mundane identification in her Transformed self's possession, the character may make a roll to uncover her identity there and then. Usually this is a Contested Wits + Empathy vs Manipulation + Subterfuge roll as the Princess makes her excuses. 

\begin{myindentpar}{10pt}
    \emph{The investigator rolls Dramatic Failure:} The investigator dismisses his evidence entirely, and takes the Incredulous Condition making future attempts to identify that Princess even harder. 
    
    \emph{The Princess wins the roll:} The investigator fails penetrate the Princess' dual identity. But gains a Clue that will help in future investigations. 
    
    \emph{The investigator wins the roll:} The investigator he identifies the Princess correctly.
    
    \emph{The investigator wins with an Exceptional Success:} The investigator also takes the Informed Condition regarding the Princess. 
\end{myindentpar}

If the Princess rolls an Exceptional Success downgrade her rivals result by one step: Exceptional Successes becomes Successes, Successes becomes Failures and Failures become Dramatic Failures. On the other hand if the Princess rolls a Dramatic Failure then upgrade the witness' result by one step.

If the investigator doesn't have a smoking gun the only other way to uncover the Princess' identity is to undertake a full investigation which requires at least five Clues to uncover the truth. More if the Princess has the Veiling Merit. A Princess should be careful because just as her enemies might investigate her Transformed Identity in the hope of finding out her secret, worried friends, truant officers, or a Princess' mundane enemies might instigate to find out where her mundane identity keeps disappearing off too.

If the players are investigating a Noble antagonist they can use the rules for investigations as written (Chronicle of Darkness p79). Unless they are distracted or dissuaded from continuing they will eventually succeed. The question is how long it takes and what dramatic complications they face on the way and as a result. But if the Princesses are the ones being investigated some tweaks will be needed: For player characters every Clue should be directly linked to a mistake they made on screen. If the Princess sneaks away to the bathroom just before her Transformed identity enters the room one too many times that could create a Clue. If she makes a call on her cell phone from somewhere her mundane identity has absolutely no business being but her Transformed identity is known to have visited that's a Clue. If she leaves some of her mundane identification behind while Transformed that is definitely a Clue.

Unless the players have gotten extremely careless Storyteller's should avoid having a supernatural antagonist follow the Princess back to her mundane life. Instead as the Princesses create more Clues you should show signs of the investigation progressing: People that the Princess helped in the past warn her about people asking suspicious questions. A surprising amount of vans are parked near areas the Princess is known to frequent in her Transformed identity. The Princess' mundane identity might even be interviewed as a potential witness. If nothing else makes sense one of the players could have a prophetic dream warning her that she's under investigation. This gives the players a chance to cover up their tracks, or if they can't remove the clues, at least do something to throw doubt upon them. 

A careful balance must be struck. Players should not be able to play fast and loose with their secrecy identity and then cover everything up at the last moment. But at the same time friends, family, and having a mundane life are thematically important to Princess, so if the friends or connections of a Princess' mundane identity are harmed the players should \textit{always} think \textquotedblleft this is fair, this happened because I screwed up.\textquotedblright\, That said if the players befriend a someone with supernatural traits in their mundane identity then that character will naturally remain exposed to the occult world and it's dangers. Even beacons, who cannot hide their Light like an unTransformed Princess, can find themselves targeted by the Darkness. 

Storytellers could consider tracking major blunders like talking about your school while Transformed separately from minor events like regularly sneaking off to the bathroom to Transform, and using the first for serious threats and the latter as a justification to introduce low threat antagonists like snooping mortal Youtube journalists for players to bamboozle or persuade. For serious threats a storyteller might consider withholding second chances if a player repeats a mistake that already caused them trouble, reducing the amount of breathing room players have to cover up old mistakes as the story progresses, or being especially harsh on mistakes made during an active investigation. 

Supernatural powers that could reveal a Princess' identity are also blocked by the innate protections guarding a Noble's Dual Identities. Introducing an antagonist who's powers could discover and attack the Princess' mundane family without any mistake on her part will hardly qualify as a fair game. Any supernatural power that would outright reveal the Princess' identity will fail: Reading a Princess' mind will not reveal her secret or directly uncover Clues, though it might provide leads to investigate. Abilities that detect the presence of the supernatural, such as the Charm Reflected Light, will not detect a Princess in her mundane identity (which is unsurprising, without an external source of power like a Bequest or a Queen an unTransformed Princess only has her Echo, which is barely magical at all). Even abilities that specifically detect the presence of Light fail. An unTransformed Princess converts a little of her Light into her Echo and hides the rest safely in her Transformed self until it is time to 
use it. The storyteller may consider relaxing this protection if the players are investigating a Princess, but not playing them. 

However supernatural abilities can certainly assist in an investigation. Reading the minds of witnesses is more effective than interviewing them. Money the Princess handled in both identities could reveal a Clue under Psychometry. If a Princess covers her house with magical defenses and someone else can sense that magic that is probably enough to earn a roll to identify the Princess there and then. An ability that detects the active use of magic could catch an unTransformed Princess using Practical Magic. 

\subsection*{Dual Identity and Intimacy}

When targeting magic against a Princess Intimacy or similar traits such as Sympathy or True Names that connect to one of the Princess' forms does not connect to her in the other form. Even if the Princess uses the same name in both identities, there's a lot of people called John Smith and they all have a unique True Name. However if a Princess interacts with the same person in both forms her Intimacy can quickly reach parity in both identities, after all her feelings are (usually) similar in both identities. Even for those who know her secret, a Princess is still two versions of the same person rather than one person who uses a disguise. If a power is currently targeting a Princess through a Intimacy connection, Transforming automatically causes it to loose track of her. 
\begin{comment}
\subsubsection*{Indiscretion \textemdash\, an optional rule}

The Dual Identity rules above require the Storyteller to keep track on which Clues the players have left behind, which antagonists are investigating the players, and which Clues the antagonists are capable of discovering. Storytellers who prefer may use this simpler system. 

The first time in a scene that a Princess does something that risks connecting her mundane identity to her Transformed identity she must roll an appropriate dice pool such as Wits + Subterfuge if she let details slip in conversation or Dexterity + Stealth to see if she avoided the security cameras when she slipped up.  
\end{comment}

\section*{Practical Magic}

In addition to the immense power each Princess wields in her Transformed identity, the Nobility can also wield Practical Magic. A subtle magic that she can call upon in both her Identities. Which abilities a Princess possesses depends upon which Queen she follows and the details can be found under the relevant Queens. 

\section*{Modified Merits}
\phantomsection
\label{sec:Merits}

\subsection*{Mentor (\textbullet - \textbullet\textbullet\textbullet\textbullet\textbullet)}

The Mentor Merit detailed in the \textbf{Chronicles of Darkness} focuses on relations within mortal or supernatural society. While both of those can be of great help to a Princess, a further option is available to a lucky few Princesses: Personal mentorship from their Queen. Even though a Queen cannot leave the Dreamlands, she still has such a wealth of power and experience that she would be far beyond a five dot Mentor. However the Queens have many obligations that demand their time besides tutoring a Princess. Therefore when applied to a Queen the Mentor Merit represents not her skill (rather than three skills, Queens can be assumed to have every skill. Even the Queen of Clubs has the raw intellect and magic to master computer programming in an hour or so if she was ever given a reason), but how much of her attention she will bestow upon the character.
%Editor: I assumed the World of Darkness rule book mentioned here is actually Chronicles of Darkness/New World of Darkness. Have changed to be so. 

As a general rule you may assume that for each Dot in Mentor you get at least one chance to talk with your Queen per Session. A chat can be just a few moments before the throne, but at high dots it can be significantly longer. After a consulting with your Queen you either gain the Advised condition or more direct clues about what you should be doing next. 

When portraying the Queens the Storyteller should always remember that they have had no contact with earth for millennia and are consequently out of touch. Therefore no matter how many dots a Princess has in Mentor and no matter how important it is; the Queen of Hearts cannot reveal to her followers the innermost secrets of the Chinese political system. Instead she can teach the Princess how to discover such secrets herself.

A Princess must follow a Queen before she can place dots in Mentor in a Queen. Neither Storms nor Mirrors will serve as Mentors. A follower of The Queen of Tears cannot place more than three dots of Mentor into her Queen. The Queen of Tear's kingdom is under siege and conforms to the normal rules of time, she just does not have enough time to spare.

Zero dots: As a Princess you are always welcome to attend court, but unless you have important business your Queen will focus her attentions on those who do. When the players have an emergency they can usually get a chance to explain and some advice. The amount of time received depends on how serious the emergency is, and whether the Princess should be able to handle it or needs significant help. Remember that both of these are opinions and two Queens might see the same situation differently. The Queen of Swords considers loosing a boyfriend as roughly equivalent to a small crisis. Even her own followers often think that's taking things too far.

With one dot a Princess can usually get the occasional snatch of time whenever there is a lull in court proceedings. Enough time to talk about recent events or ask questions and get a brief bit of advice.

At three dots the Queen often has your name in the royal itinerary for time with you alone, or perhaps group sessions with your Nakama. You have the opportunity to talk at length about your progress in the war of hope and receive private tutoring. You also know what your queen is like when she has an opportunity to take off her crown, so to speak. Generally speaking a Princess with three dots of Mentor has something special about her that caught her Queen's eye. Something the Queen is especially keen to nurture. She may be the only Onceborn in the Court or have the largest following among the Queen's court. Perhaps she was someone the Queen respected in a past life or maybe the Princess works full time as one of the Queen's agents upon Earth and must report in often (This can be especially helpful for beginner players as it both justifies clear goals for the character and advice from the Storyteller).

At five dots the Princess is spending at least an hour or so with the Queen most nights, and if they don't show up the Queen may even get worried and send one of her other followers to check on the Princess. This level of attention is exceptionally rare and Characters who have five dots are either vital to the Queen's agenda or have an exceptionally strong personal connection. Perhaps the Queen believes the Princess might become a Queen herself and has dedicated herself to making sure it happens. The Princess may be the first Ambassador to a concept the Queen cares for greatly. Alternatively the Queen and Princess might have an immense personal bond of friendship, or maybe the character is the Queen's consort (though this is usually only an option for followers of Swords).

\subsection*{Retainer (\textbullet+)}

While Princesses have as much benefit from a normal human retainer as anyone else they may also seek out Beacons, accept the fealty of Sworn, or travel with Shikigami. 

For a Beacon or Sworn, their magical abilities are basic add a single additional Dot to their effectiveness, and merit cost, as a Retainer, this remains true even when one person wields both the magic of a Beacon and a Sworn. 

Shikigami who work side by side with or for a Princess can also be represented by the Retainer Merit. Since the typical Shikigami is a person of note in the Dreaming Kingdoms and has access to magic they're rarely worth less than four dots as a retainer. As a special rule a Nakama can split the cost of a Shikigami between them, after all one of a Shikigami's most useful abilities is serve as a magical relay between all the Princesses it is bonded with. Shikigami who come and go as they please or who laze around the Princess' home and dispense advice or clues to drive the plot forward in-between binging on cakes and MMORPGs do not need to be represented by a Merit. 

\subsection*{Sympathetic (\textbullet)}

Princesses are good, perhaps too good, at letting others get close. The Sympathetic Merit can be brought by Princesses for only one dot. 

\subsection*{Taste (\textbullet)}

A Princess' Transformed self, and the Regalia she wears are the innermost expression of who she wishes to be. This expression is influenced by her culture, this is why Princesses in Europe who wish to be \textquotedblleft pure\textquotedblright\, often have white as a prominent motif in their Regalia, while Princesses in much of Asia often wear white if they wish to work with ghosts. When used to study Regalia the usual questions answered the Taste Merit can be asked about the Princess herself. On an Exceptional Success gain the Informed Condition, applying to the Princess' self identity. 

\textbf{Drawback:} Studying a person isn't quite the same as studying an object, the Princess can contest the roll with Manipulation + Subterfuge to cloud the issue. If the Princess disguises her regalia with magic the benefits of Taste are no greater than when studying the fashion of any other disguise.

\subsection*{Virtuous}

Being virtuous is a prerequisite for Blossoming, consequentially a second Virtue is part of the Princess template without requiring a Merit and thus a Princess may not take the Virtuous Merit.


\section*{Noble Merits} 
The Hopeful may buy Merits with Attribute or Skill prerequisites that they only meet when transformed. If they do so, they only gain the benefit of the Merit in question while in their transformed identity. You may convert transformed dots to mundane dots unlocks Merits with prerequisites in the mundane identity if the character's new mundane traits are high enough to support the Merit.

Below are new Merits related to the Princesses. Unless otherwise noted, all these Merits are reserved to the Hopeful. 

\subsection*{Circle (\textbullet\nobreak\hspace{.16667em}-\nobreak\hspace{.16667em}\textbullet\textbullet\textbullet\textbullet\textbullet)}

The human mind was not designed for large numbers. When the suffering of millions feels incomprehensible, you can turn to the small everyday problems and triumphs of friends and family to remind yourself what you are fighting for. The Circle Merit represents healthy relationships between a Princess and her family, close friends, or mentors, who are ordinary mortals (at most they may have supernatural merits like Psychokinesis, Beacons qualify too.) The more dots a Princess has, the more or the stronger these relationships are and the healthier her social life is. Describe the people your character is intimate with when you buy dots in Circle.

Once per day a Princess with this Merit can spends an hour socializing with her Circle and regain up to (Circle * 2) - Shadows Wisps. She doesn’t have to speak to every friend she has in an hour, like most people she probably sees different friends or family members at different times. 

Friends and family are good for more than Wisps, they also provide a support network which provides pervasive modifiers to Compromise rolls. A Princess with two dots of Circle gets +1 to all Compromise rolls, at four dots the bonus increases to +2, however a Princess with no Circle has a -1 penalty to all Compromise rolls.

\textbf{Drawback:} Friendships require investment, family is more forgiving but even a family relationship can grow distant. If a Princess goes a significant amount of time without socializing with her Circle she may lose dots. 

\begin{mybox}     

\subsection*{Transformed Merits}

Sometimes a Princess has a merit which only applies in her Transformed identity. Her perfected self may be significantly prettier than her mundane self (that's a common one), she may transform into a tougher person. Alternatively a Princess may live off the grid in her mundane identity but have achieved fame for her exploits in her Transformed identity. Every time a player buys a Merit she can declare it only applies to one identity and take a Beat. Transformed Attributes and Skills count towards any Prerequisites, but only in the Transformed identity of course. 

If the Princess puts in some effort most merits can be expanded to cover both identities. A Princess who is Fleet of Foot in her Transformed body could train herself at running in her mundane body. This costs nothing more than repaying any Beats given for taking a Transformed Merit. 

This only applies to merits that are part of the character. Resources, Contacts, and any merit that represents something external to the Princess do not provide Beats. They are assumed to belong to one identity by default, unless the Princess wants her Contacts to know about both her identities.

\end{mybox}

\subsection*{Emotional Intuition (\textbullet\textbullet\textbullet)}

People are always being told to follow their heart, and when your emotions are the source of magical powers following your heart can lead you right. Once per chapter you may ask the storyteller one of the following questions. Roll Belief, on a Success the storyteller answers truthfully. On a failure you get no answer. An Exceptional Success lets you ask a second question. If a character has been attempting to deceive the Princess they get a roll to contest any questions asked about them. Even if they are not present they can contest, the defensive roll represents any false impressions the Princess retains from past deceptions. The Storyteller should roll the target's Manipulation + Subterfuge in secret. If the target wins the answer is what they would want the Princess to think.

\begin{itemize}
 \item How will this character react to this action?
 \item How does this character feel about this thing?
 \item How does this character feel about that person?
 \item What is this character hoping to achieve? 
 \item Can I trust this character? 
\end{itemize}

\subsection*{Iaijutsu Kensai (\textbullet\textbullet)}

You have learned to use Transformation itself as a weapon, manifesting your Regalia mid-motion to catch your foes by surprise. Whenever your Kensai weapon is not manifested and you Transform reflexively you may roll Wits + Brawl or Weaponry vs Wits + Brawl or Weaponry. If you win you may make an attack using your Kensai weapon that bypasses your opponent's Defense. 

A character who witnesses you using Iaijutsu Kensai is immune to it for the remainder of the scene. 

\subsection*{Onceborn (\textbullet\textbullet)}

Available at character creation only.

A reincarnated Princess is sure to Blossom sooner or later, and usually sooner. For anyone else the requirements are much stricter and yet you qualified. You were born with a noble heart, you lived a life that pushed you to achieve ever greater heights, and you reached the Light itself. 

Your lifetime of experience means your Belief begins at 8, not 7.\\

\noindent\textbf{Drawback:} You have no experience from past lives to call upon. 

\subsection*{Palace (\textbullet)}

There is a difference between a house and a home. A Palace is not a place to live, though a Princess may live in one. A Palace is a place with such an intimate emotional connection to a Princess that her magic works better there than anywhere else. A Palace automatically has the Consecrated Condition, which lasts indefinitely just so long as the Princess' emotional connection remains strong. Any additional benefits of a Palace may be represented by Merits such as Safe Place or Gallery, or as Equipment. 

\subsubsection*{Palace Veiling (\textbullet\,\nobreak\hspace{.16667em}to\nobreak\hspace{.16667em}\textbullet\textbullet\textbullet)}

Something odd happens to some Palaces, when its mistress deliberately associates her domicile with her mundane self and avoids associating it with her Transformed self. The effects of her Dual Identity seeps into her home, preventing others from connecting it to her Transformed Identity. Each dot assigned to Palace Veiling provides a -1 penalty to any attempt to find the location of the Palace when searching for the home of it's mistresses' Transformed Identity. Searching for her Mundane Identity's home has no penalty. Anyone inside the home suffers the same penalty to notice any evidence of the Princess' identity. This is not infallible, a -3 penalty is not enough for her to greet guests in full Regalia and though she might get away with leaving her magic wand on her desk, don't risk it. 

Charms that take advantage of the Palace's innate Consecration enjoy an even greater protection as though they were protected by the Grey Masquerade Charm. Passive supernatural abilities such as the Charm Reflected Light automatically fail while active abilities such as the Charm A Magical Girls Eyes must beat you in a Clash of Wills to notice any Charms that take advantage of your Palace's Consecrated Condition. You add your dots in Palace Veiling to the Clash of Wills. 

Finally a Veiling also serves to protect those inside from hostile scrying, penalizing all attempts to scry into the affected area by the dots in this Merit. This does not stack with any other Merits that penalize scrying \textendash\, apply the larger penalty if a character with such a Merit scrys upon a veiled Palace.

Palace Veiling has no effect on anyone who already knows both the Princess' two identities and has connected them to each other.

\subsection*{Royal Tongue (\textbullet)}

Prerequisite: Sworn or Hopeful. It might be possible for other supernatural creatures to find a way to mimic its effects, but that would require the active use of powers (such as the Mind Arcanum or a Majesty Devotion), and so they cannot learn the Merit.\\

\noindent You know how to speak the Royal Tongue. You can use it to enhance your Charms.

As languages go the Royal Tongue is unique, it is quite possibly the densest language there is. Meaning is conveyed not just by words but also tone and the placement of stresses, two different words could be completely identical except for where it's grammatically acceptable to use them. A monosyllabic prefix or suffix can completely change a word within it's conceptual space; the Royal Tongue has over a thousand words for friend, (and just as many for lover or sister or brother) each defines the exact details of the relationship and not only do they all come from the same root word, they all sound like that root word.

The hardest thing to do with the Royal Tongue is speak a straight sentence. Between the grammatical rules and the enormous amount of information in every word a simple sentence like \textquotedblleft thank you for helping\textquotedblright\, would, if translated into English, look like two paragraphs of flowery purple prose about the inherent virtue of teamwork. When asked how such a language came about the Queen of Hearts implied that it was deliberate, the Queen of Spades burst into laughter and the Queen of Diamonds (who's fluent in, and prefers, most Earthly languages) gave a put upon sigh.

Mortals can hear the the Royal Tongue but they cannot understand it. Many Princesses swear blind that the Royal Tongue uses magic as well as sound waves as a carrier to pack even more information into each word. The cost of this is that only Princesses, Sworn and Beacons can learn to understand the Royal Tongue. Beacons can't even speak it. To anyone else it sounds like someone singing in an unknown melodic language, the emotional intent comes across clearly but the words are incomprehensible. The Royal Tongue can't be recorded by mundane means either, a Princess who hears a recording of the Royal Tongue catches at best the occasional word, it sounds as though someone has removed 40\% of the syllables and all the grammar. This is considerably more confusing than it would be in a normal language.

The most useful aspect of the Royal Tongue is to enhance a Princess' Charms (Though some Troubadours insist it's use in lyrics and poetry is more useful... Deep purple prose is an acquired taste). If you spend a turn speaking the Royal Tongue, declaring that you will use a Charm on a target, and forcing the qualia, the experience and feeling of what you are about to do, into the musical speech, your activation roll for the power gains the 9-again quality. 

Use of the Royal Tongue is a supernatural ability, and triggers the sense of individuals with Unseen Senses related to Princesses or the Light.

\subsection*{Second Calling (\textbullet\textbullet)}

Available at character creation only.\\

\noindent Your character feels two Callings in her Heart. Choose a second Calling for her; she may recover Wisps from that Calling just as well as first, and she may follow Dreams from either Calling. 

She does not gain an extra transformed Attribute dot or affinity to any Charms from her second Calling.\\

\noindent \textbf{Drawback:} Your character is bound by the Oaths of her second Calling, just as much as those of her first - both sets of Oaths are breaking points for her.

\subsection*{Tomoyo's Touch (\textbullet)}

Your character is unusually adept at manipulating her Regalia. When a Charm that creates Regalia specifies that it's parameters may be modified by a Transformation action your character can make those modifications, as well as switch which of that Charm's Upgrades are active, as a Reflexive action without needing to spend a Wisp or make a Transformation roll. She may do this once a turn, but when she does so she may affect as many pieces of regalia as she wishes.

\subsection*{Troupe Magic (\textbullet)}

\textbf{Prerequisite:} A Commonality of Dedicated. \\

\noindent Your character belongs to a performing Nakama who've learned to reinforce each other’s magic, as well as their vocals. You may perform Inspire Charms as a Teamwork action even with secondary actors who do not know the Charm providing the secondary actors are part of your Nakama. 

\subsection*{Unfinished Business (\textbullet\nobreak\hspace{.16667em}-\nobreak\hspace{.16667em}\textbullet\textbullet\textbullet)}

Something in your character's past life still haunts her, and when an opportunity to address the issue arises (which usually means tracking down the reincarnation of someone else involved) your character's past self may replace her current personality for a short while. Once per Session, while you are Transformed, you may elect to have your character's past self take over and peruse her goal. You get two experience points per dot of this Merit to represent the skills and abilities your presumably older and more experienced incarnation knew at the moment she died, you do not have to spend the experience the same way each time. If your past self causes a significant hardship for your present self, take a Beat.

Princesses who manifest their past self this way don't manifest all of their former self. In behavior and attitude they are identical, but they are clearly people outside their own time. They have trouble forming emotional connections to present day people or situations, which is extremely unusual for a Noble, and relate to the world as though they were just reading about it in a particularly dry history text book. They show a similar lack of concern about the fact they are dead. Only something which relates to her unfinished business or her own time can form an strong emotional reaction.

Not every Princess agrees with her past self and others simply don't have the time. Some find closure through therapy rather than a quest, and others seek inner peace through meditation or in their Crawlspace.

It is recommended that Characters with Unfinished Business also take the White Rabbit Dreams Merit as that ability makes a great way to provide clues to her Unfinished Business and find situations where you can bring out your character's past self.

If you resolve your unfinished business remove the Merit and refund it's cost. 

\subsection*{Veiling (\textbullet\nobreak\hspace{.16667em}-\nobreak\hspace{.16667em}\textbullet\textbullet\textbullet)}

Some Hopeful's perfected selves differ so much from their mundane selves that they are excessively difficult to identify, even beyond the norm. A Princess with Veiling is one such; whenever an observer tries to connect her two identities, they subtract Veiling from all rolls, and each dot of Veiling adds an additional Clue as a requirement to uncover her identity. 

If a Princess has Veiling and Anonymity the two merits have distinct roles. Anonymity makes it hard to find information about the Princess (or one of her identities). Once you have that information Veiling makes it hard to connect it to her second identity, or even realise she has two identities. 


\section*{Common Merits}
The following merits are available to both Princesses and mortals. 

\subsection*{Bequest (\textbullet +)}

\textbf{Prerequisite:} Princess, Beacon, Shikigami, or Sworn

Your character possesses an object that contains a Charm bound within. This might have come from within your own soul, to be granted to another or used to boost your own powers, or might have been obtained on a quest to the Dreamlands. Some senior members of the Nobility (such that there are) make a habit of gifting trusted friends with them for, despite the cost, they can give a Princess access to powers that she might not otherwise be able to reach.

Every Bequest, like the Princesses, has two appearances: a mundane one from the vessel, and a glorious one from the Charm. And, like the Princesses, a Bequest's power can be used only when the Bequest is transformed. Transforming a Bequest is somewhat more difficult than transforming yourself; you can spend a Willpower point to do it reflexively (not a Wisp) or make a transformation roll of your Belief (not Belief + Inner Light.) A Bequest, once transformed, remains so for 1 scene. You need not be transformed to transform or use a Bequest, and you do not have to transform a Bequest yourself to use it.

Any of the Hopeful, whether transformed or not, may recognize a Bequest as such; she rolls her Inner Light the first time she touches it in a scene. Bequests also show up on an Unseen Sense for the Light.

Bequests based on permanent Charms apply their Charm's effect to the person holding them as long as they are transformed. Bequests based on activated Charms require the person holding them to pay the Charm's cost, and then apply the Charm's effect. If a Bequest is based on an upgraded Charm, every upgrade included in the Charm applies every time the Bequest is used, and you must pay the upgrades' cost in full \textendash\, a Bequest, unlike a Charm, is not flexible. If you have an Invocation that's compatible with the Bequest's Charm, you may apply it when using the Bequest, with the same effect as it would have on the Charm. If the contained Charm's effects very by Inner Light, the user's Inner Light applies as normal.

Some Bequests simply cannot be used by anyone without an Inner Light of their own (that is, Princesses and Queens); This includes Bequests that contain the Charms: Bequeath, Accept Fealty or Long May She Reign. Bequests containing an Avatar Charm cannot be used by anyone who does not have four dots in the related Invocation.

Many Bequests carry a pool of Wisps within themselves, which their holders can spend to activate them (but not for any other purpose.) The capacity of the pool is 10 Wisps. To add Wisps to the pool, the Bequest’s owner has to carry out a duty while the Bequest is on her person, roughly equivalent to a duty of a Calling (and using the rules for them); the specific duty, and the dice pool the owner rolls, are fixed at the Bequest’s making. A Princess cannot extract Wisps from a Bequest, without the Charm Charge and the upgrade Accepted. If the Bequest itself provides the Charm charge it can transfer it's own wisps freely, as well as wisps from the user's pool.

A Bequest's base Merit dot cost is equal to the level of the Charm (so, for example, Charge, a Govern 2 Charm, costs two dots), plus one dot per upgrade, plus one dot if the Bequest has a Wisp pool. For more details on Bequest creation, see the rules for the Bequeath charm.

\subsection*{Destiny (\textbullet\nobreak\hspace{.16667em}-\nobreak\hspace{.16667em}\textbullet\textbullet\textbullet\textbullet\textbullet)}

Your character’s fate is a thread of gold in a tapestry of cloth. She may be fated to shake nations or just change her own family but in the cosmic scheme of things she \textit{matters}, perhaps more than she will ever know, and the universe itself guides her actions down a fated path. You gain one point of Destiny for every dot in this Merit, spent points refresh at the start of every chapter. A point of Destiny can be spent before rolling to turn any single mundane roll into a Rote Action, or to reroll a single mundane action after you see the result of the roll, the second result stands. 

When taking this Merit name your character's destiny, the storyteller may veto anything that doesn't fit the story. If your character successfully achieves her Destiny remove and refund this merit.\\

\noindent \textbf{Drawback:} Your character's success is not guaranteed. Every fate has a escape clause, a way that it might be prevented. This is the destiny's Doom. A person, object, situation or even a concept that can drag your character into a horrible fate. It either kills her, or leaves her alive and suffering. If the Princess doesn't like her destiny the Doom should be even worse. (See the Fate optional rule [CofD 177] for examples of Dooms.) 

Mechanically if your Doom is active in the current scene you cannot spend points of Destiny. Beyond that it is up to the Storyteller to ensure that conflicts involving the Doom are climatic moments with high stakes for failure.

\subsection*{Entwined Destiny (\textbullet\textbullet\textbullet)}

You were bound to another before you were born, and both events and your inmost will urge you towards a specific role in his life. When you take this Merit, choose the nature of your relationship with its subject -- true love, rivalry, protection, loyalty, and enmity are all possible, and if none of these fit you may define your own, with the Storyteller's approval. You may also choose the subject's name and general nature, or leave that up to the Storyteller to decide. Once a scene, whenever you succeed in a significant task that plays into your destined role, you regain 1 point of Willpower. While there can be exceptions, nearly all occurrences of an Entwined Destiny are mutual, if you are destined to be somebody's rival they share an equal destiny to be your rival. 

Entwined Destiny cannot directly influence an person's choices, but when you have this merit Fate will attempt to nudge you in certain ways. For example if you have an Entwined Destiny of love Fate cannot make you feel love. Instead Fate can and will create coincidences where you and your destined bump into each other in circumstances appropriate for heart to heart talks. If the destiny is natural it can also be assumed that your destined true love are compatible people, your destined ward will need a protector, or similar.

\textbf{Drawback:} In any scene when you have an opportunity to fulfill your destined role in dramatic action, and deliberately forgo it, you take a -2 penalty on all your actions until the scene ends. For example, if your destined rival challenges you, you must accept; if you discover your destined enemy's plot, you must oppose it; if your destined ward is in peril, you must rescue him.

\subsection*{Gallery (\textbullet\nobreak\hspace{.16667em}-\nobreak\hspace{.16667em}\textbullet\textbullet\textbullet)}

Humans are social creatures and Princesses even more so. A Gallery is a set up designed to help with anything from dinner parties to theater shows. Each Gallery covers one Specialty. When performing any Expression, Persuasion or Socialize action which fits under a Gallery's specialty the Princess gets a bonus equal to the dots, so long as she inside the Gallery at the time. 

\textit{Example. Princess Gwendolyn of Hearts wants the mayor to invest more in the town's cultural heritage. Gwendolyn has Gallery (\textbullet\textbullet\, Dinner Parties) and has invited the mayor to discuss her proposals over a meal. She gains +2 to all Persuasion rolls during the evening.}

\subsection*{Locus of Destiny (\textbullet\textbullet\textbullet\textbullet\textbullet)}

\textbf{Prerequisite:} Destiny. 

Your character's destiny is so potent that fate will not leave it to you alone. For every dot in Destiny you gain one free instance of the Entwined Destiny Merit, if you gain more dots of Destiny new companions or rivals will soon show up. The player and the Storyteller should work together to decide how the various other characters fit into your character's destiny. Your destined mentor or mentors could train you for the task ahead, while in a cruel twist of fate your destined rival could even be your destiny's Doom; forcing your character to overcome your fate in order to fulfill it. 

\textbf{Drawback:} Your Destiny is especially hard to resist and with fate offering temptations even your closest companions can be surprisingly pushy or even unsympathetic if you act against your fate. At any time your character is trying to avoid their fate all characters who you have an Entwined Destiny with switch their destined role to \textquotedblleft getting you back on track\textquotedblright. Though they will do so in a manner that befits their personality, and thus probably their original role. When the character is once again moving towards their destiny everyone's destined role returns to normal. As always Fate cannot actually force your character's companions to turn upon you, but there is a world full of potential destined companions and Fate has a habit of picking ones who'll stick to the path.

Your character's destined companions will not intentionally push you towards your Fate's Doom. However their attempts to push you towards the destiny's conclusion may throw you into the path of the Doom. Their role is to keep you on track, regardless of whether your doom is planted right in the middle of the tracks or miles off to the side. 


\subsection*{Lucid Dreamer (\textbullet\textbullet)}

Your character is aware of when she is dreaming and has the ability to act consciously within her own dreams. This is especially useful if someone else is inside your dreams, as it gives you the ability to interact with and oppose intruders. While in their own Dreams a mortal can roll an appropriate Mental Attribute + Integrity to perform any physical action imaginable. Whether that's imagining a gun and shooting an intruder with it or imagining the entire dream flooded with lava. A success is a success, dream guns and lava are equally powerful barring special circumstances. 

Lucid Dreaming is common among Nobility, and even more frequent among their Transformed Selves. Maintaining awareness in the Dreamlands, a universal ability for Princesses, is not much different from remaining aware in a regular dream. When only the Transformed self is a lucid dreamer it is the Princess' Transformation within the dream, not her physical body, which applies. This means that a Princess who is attacked in her dream and instinctively transforms achieves lucidity. 

While asleep a lucid dreaming Nobel uses Belief instead of Integrity, though most use their Charms instead as they are often significantly more powerful.

\subsection*{Mandate (\textbullet\nobreak\hspace{.16667em}-\nobreak\hspace{.16667em}\textbullet\textbullet\textbullet)}

\textbf{Prerequisite:} Princess or possession of a Bequest with a duty. \\

\noindent Every Calling has it's Sacred Oaths, the means by which a Princess feels is her duty to help the world. The Mandate Merit represents not an emotional connection to one's Calling, but a practical ability to fulfill it. Any Grace could take it upon herself to help the local bully apologize and make amends but she'll only come across so many bullies in a month. A Grace who works as a relationship councilor, she'll be seeing clients every day. A Champion could work on the police force, a Troubadour could have a popular blog where she posts poetry. 

Whenever a Princess rolls to regain Wisps through Voluntary Work a successful roll grants an additional Wisp per dot of Mandate.

\textbf{Drawback:} A lot of Mandates are responsibilities as well as opportunities. A Mender who's parents own an animal rescue centre can leave the animals in her parent's hands while she goes of to deal with some Darkspawn. A Mender working as a Doctor at the local hospital will be in trouble if she doesn’t put in her hours. Generally Mandates without responsibilities tend to be worth few dots. 

The Storyteller is encouraged to see a Mandate as a way to start interesting plots, not just a time sink that keeps players from the adventure. Darkspawn victims could come into the hospital emergency ward. A police officer could assigned and interesting case, mundane or supernatural. A blogger might meet a fellow blogger who reports on supernatural events. 

\subsection*{Regal Bearing (\textbullet\nobreak\hspace{.16667em}to\nobreak\hspace{.16667em}\textbullet\textbullet\textbullet\textbullet\textbullet,\nobreak\hspace{.16667em}Style)}

\textbf{Prerequisites:} Composure 3 Socialize 3.

Your character carries herself with an air of sophistication and grace. Whether she rubs shoulders with European royalty or commands respect by walking through the door of an inner city church in Sunday best there's no denying she is noble.

\textbf{Figure of Authority (\textbullet)} Your character was born to authority and wields it well, whenever she adds Status to the dice pool of a mundane action she gains an additional die. 

\textbf{Effortless Elegance (\textbullet\textbullet)} Your character knows every unwritten rule of manners and never misses a beat. In a formal situation Socialise specialties also provide their full bonus to Composure rolls within the specialty's domain, so long at the roll is mundane and not supernatural. If the roll is Composure + Socialize the Specialty applies twice. 

\textbf{Stiff Upper Lip (\textbullet\textbullet\textbullet)} A gentleman is not a fancy suit and elocution lessons, a gentleman is a way of life and your character lives this truth. At any time you can spend a point of Willpower to lesson the effects of a Condition representing a mundane mental or emotional effect upon your character. For the remainder of the scene any dice penalties that Condition applies to Social rolls is reduced by Composure. This Merit cannot affect Conditions representing a supernatural or physical effect (such as wound penalties or intoxication).

\textbf{Dry Wit (\textbullet\textbullet\textbullet\textbullet)} Whether your character looks down her nose at the help or reminds the room that an old family name does not make a pompous buffoon into a gentleman; she does with such grace and subtlety that nobody can accuse her of impropriety without desecrating their own dignity. Any attempt to tarnish your character's reputation based on her words, but not her deeds, takes a penalty equal to her Expression. This does not apply when her words provide evidence of her deeds. 

\textbf{Manners Maketh Man (\textbullet\textbullet\textbullet\textbullet\textbullet)} Your character radiates authority and poise, she is a pillar of strength in times of chaos. Whenever your character begins Social Maneuvering with a character who is not opposed to them and currently has a non-Persistent Condition representing a mundane mental or emotional effect they automatically open a number of doors equal to half your character's Composure, rounding up. This Merit provides no benefit when the target is suffering from Conditions representing supernatural or physical effect (such as wound penalties or intoxication).

\subsection*{Populist Rhetoric (\textbullet\nobreak\hspace{.16667em}to\nobreak\hspace{.16667em}\textbullet\textbullet\textbullet\textbullet\textbullet,\nobreak\hspace{.16667em}Style)}

\textbf{Prerequisites:} Presence 3 Manipulation 3.

Your character has been trained in making friends and influencing people. She might be a politician or a cult leader, even if she is not she sees words as tools to make people do what she wills. Such power can be used for both good or ill.

\textbf{Anecdote (\textbullet)} Your character brings up a heart-wrenching story. Arguing against her is like arguing against little orphan Annie. Roll Presence + Expression vs Composure + Expression, on a Success your target has a -2 to social rolls against you on this topic for the rest of the scene.

\textbf{Polispeak (\textbullet\textbullet)} Your character can talk at length while saying nothing. By spouting a never ending stream of ambiguous statements and platitudes you can dodge questions while appearing to give answers. When someone is debating against you, you may reflexively forgo your action for a turn to subtract twice your manipulation from your opponents Persuasion or Expression rolls. This works best at stalling for time or when you can't win a debate and hope to simply minimize the scale of your loss.

\textbf{Catch phrases (\textbullet\textbullet\textbullet)} People are naturally tribal and like a social chameleon your character can fit in anywhere by spouting catchphrases and speaking with the correct lingo. Once per scene you may spend a point of Willpower and gain a bonus equal to half her Subterfuge to all social rolls with a certain group. This cannot be done if the target group has Prerequisites your character does not meet, an adult could not join a teenager only group for example. Using this technique is not declaring loyalty to a group, in fact it's probably best to avoid doing so, but if she's good people will naturally feel more comfortable around her.

\textbf{Drawback:} By getting in with a group your character is placing herself outside other groups. She gains an equal penalty to those who dislike the group, for example if she creates a bonus with Republicans she'll take a penalty while talking to Democrats for the rest of the scene. 

\textbf{Encouragement (\textbullet\textbullet\textbullet\textbullet)} Your character can convince people to associate what they like about themselves with herself. You try to treat all people politely, that's not because you're a good person. That's because you're a Christian/Liberal/Humanist/Feminist/Conservative. You can apply a variant of the Inspired Condition to your listeners, it may only be Resolved on actions in line with their Virtue but when it is resolved you open a door.

\textbf{Rally (\textbullet\textbullet\textbullet\textbullet\textbullet)} Your character can call play the crowd like a fiddle, alienating and isolating her opponent. Both your character and her opponent pick a social merit that gives you influence over people (Usually: Allies, Fame, Mystery Cult Initiation, or Status though others may be aloud. The people represented by that merit must be present, be they in the audience, waving placards outside or posting comments on the blog with the debate). Roll Presence + Persuasion + Merit Dots. For every (Targets Composure + Merit dots - your successes) rolls the target must spend one point of Willpower to keep going in the face of such hostility. This does not count as the target's one willpower per turn.

\begin{myindentpar}{10pt}\textit{Example: Gwena\"{e}lle is a columnist having a public debate about obesity with Desmond, the director of public relations in a greedy fast food chain. Gwena\"{e}lle rolls Presence + Persuasion + her Status in the newspaper to fill her next column with catch phrases that sound a call to arms. She gains four successes. Desmond has a Composure of 3 and chooses to play his Status 4 in the corporation, getting subordinates to ghost write and give him some emotional distance. Every third roll Desmond loses a willpower as his corporate blog is hammered with aggressive comments by Gwena\"{e}lle's supporters. }\end{myindentpar}

\subsection*{Scientific Rhetoric (\textbullet\nobreak\hspace{.16667em}to\nobreak\hspace{.16667em}\textbullet\textbullet\textbullet\textbullet\textbullet,\nobreak\hspace{.16667em}Style)}

\textbf{Prerequisites:} Intelligence 2, Academics or Science at 3 or at 2 with a specialty in academic methodology or practices.

For your character a debate isn't about winning or convincing others, it's about spreading the truth. Scientific Rhetoric works best against the open minded, even if your character loses the debate she learns something.

\textbf{First Principles} (\textbullet): Quantum nonsense, pyramid power, \textquotedblleft I'm not a doctor but I play one on TV". People have been using scientific lingo to promote rubbish for years, to a real scientist it usually sounds ridiculous. As a reflexive action you may inflict a -1 penalty on anyone making an argument based on science or academics (the fields, not the skills). This does not work if their science is sound, assume that a specialty or three dots in a relevant skill confers immunity. Less for basic topics.

\textbf{Citation} (\textbullet\textbullet): You don't need to prove your point, just point to someone else who has already proved it for you. You may apply the Library merit to social rolls, provided you have access to your Library and your conversation partner lets you fetch books. Internet access is worth +1, unless you have a Library of bookmarks to hand.

\textbf{Just the Facts} (\textbullet\textbullet\textbullet): A quick tongue doesn't matter, for the facts speak for themselves. You may fall back to established facts, allowing you reroll one failed social roll per scene providing it is on a scientific or academic topic.

\textbf{Burden of Proof} (\textbullet\textbullet\textbullet\textbullet): You stand your ground and insist debate will be decided on the facts. Make a contested Manipulation + Persuasion roll (or a Presence + Expression roll when debating in front of an audience). If you win you and your opponent must use an appropriate mental skill (such as Science or Medicine for a debate on evolution) in place of social skills for rolls connected to that debate.

\textbf{Skepticism} (\textbullet\textbullet\textbullet\textbullet\textbullet): There is a reason reproducibility is one of the principles of the scientific method. After a social encounter has opened one of your Doors you can take the time to double check the other guy's claims. Is that salesman's product really the best in it's field? Make an extended research roll with a target equal to the dice pool used to the door. If you succeed that door closes again. You may only ever use this maneuver once for any door and it fails automatically if the claim you are investigating turns out to be true or is too subjective to be true or false.


\subsection*{Spiritual Temperance (\textbullet\nobreak\hspace{.16667em}to\nobreak\hspace{.16667em}\textbullet\textbullet\textbullet\textbullet\textbullet,\nobreak\hspace{.16667em}Style)}

\textbf{Prerequisites:} Composure 3, Academics 2 or with a specialty in an appropriate school of thought.

Your character is attuned to the spiritual side of the world. She may be a Buddhist monk, a Christian monk, or simply read a lot of paperback philosophy. Her spiritual attunement helps him above the temptations of everyday life. Spiritual Temperance excels at resisting persuasion, it is least effective when directness is a virtue and against those who cannot understand it's hidden meanings.


\textbf{Unclouded Eyes (\textbullet):} Like the fool your character favors the simple explanation. She plays little attention to distractions or complex word games. Ironically this makes her exceptionally hard to fool, all subterfuge rolls against your character take a -1 penalty.

\textbf{Cessation of desire (\textbullet\textbullet):} By focusing on the spiritual your character has freed herself from temptations. Mundane social rolls gain no benefit from tempting her with her vice.

\textbf{Drawback:} No vice means no Willpower. They player may choose for their character to indulge in their Vice but then the provider gains the usual bonuses. 

\textbf{Mediation (\textbullet\textbullet\textbullet):} Your character speaks peace to the world. So long as she gives her full attention to mediating a conversation all characters treat their impression as one level higher. As soon as your character stops the impressions return to their usual level. \textit{Example: two people with an average impression can roll to open doors once per day, providing that each roll was made while your character is mediating.}

\textbf{Drawback:} Your character cannot use mediation when she wants to open a door herself.

\textbf{Koan (\textbullet\textbullet\textbullet\textbullet):} Your charter tells a short story or gives a riddle with layers upon layers of hidden meaning and a kernel of wisdom hidden at it's core. Roll Expression + the target's Wits. On a Success the target gains the Informed Condition, however when they resolve the Condition open a door. Understanding and following your koan is another step along your path.

\textbf{Inner Peace (\textbullet\textbullet\textbullet\textbullet\textbullet):} Your character has risen above desire, she could live comfortably on nothing but rice and water. This lack of desire makes it incredibly hard to push her into anything. In any social roll where your character's only goal is to resist she may roll add Integrity (or Belief) to her dice pool.

\subsection*{Taint Awareness (\textbullet\textbullet)}
\textbf{Prerequisites:} Beacon, Sworn, Shikigami or Hopeful.\\

\noindent Your character's sensitivity to corruption from the Darkness is so finely tuned that she can perceive it from miles away. Whenever a creature of the Darkness enters our world from the Dark World, or the size or Severity of a Tainted area increases within Inner Light miles of your character (1 mile for a Beacon or Sworn) she is instantly and unambiguously aware of it. She may roll Wits + Sensitivity to acquire a sense of the direction and rough distance to the Tainted area or the place where the creature crossed over.

Some character's are especially sensitive, however these lucky or unlucky people have senses so refined that it is painful to feel an intrusion. You may, for the same Merit Dots, opt to have your character sense Taint or intrusions within Inner Light * 3 miles and gain a +3 bonus to determining the location and distance. If you do so you take the moderate version of the Sick Condition for one scene as you sense the Taint through a medium of nasuia, headaches, uncomfortable periods, cramps or some other painful medium. The exception is if you are currently in your Transformed Identity, a Princess' magical body is better able to regulate supernatural senses.

\subsection*{White Rabbit Dreams (\textbullet\nobreak\hspace{.16667em}-\nobreak\hspace{.16667em}\textbullet\textbullet\textbullet\textbullet\textbullet)}

\textbf{Prerequisites:} The ability to travel to the Dreamlands.\\

\noindent Your Character has an especially strong connection to the Dreamlands; when she sleeps she is called to witness prophecies and portents applicable to her life and her goals. When her character sleeps for long enough to reach REM sleep, the player can declare that her character witnesses a prophecy in the Dreamlands or relieves a memory that contains relevant clues to the current story. Then, ask the Storyteller one or more yes or no questions about her character’s life, surroundings, a task at hand, or the world at large. The Storyteller must answer these questions truthfully. Every Session the player may ask one question per dot in this Merit, they may use all the questions at once or split them over several nights. 

\section*{Mortal Merits}

The following merits are only available to Mortals. 

\subsection*{Dedication (\textbullet\nobreak\hspace{.16667em}-\nobreak\hspace{.16667em}\textbullet\textbullet\textbullet\textbullet\textbullet)}

Prerequisite: Sworn or Shikigami \\

\noindent A Princess' Light shines forth into the world and those Sworn to her service reflect that light becoming brighter themselves. Dedication represents one Princess who the Sworn respects and follows, or one Princess a Shikigami advises. The Merit's rating reflects the scale of that intimacy. 

When the character spends time with the Princess (or a Queen), either socializing or assisting with her duties, he may roll Empathy + Dedication once per scene or once per hour during downtime. Each success grants one Wisp, no more than two times Dedication Wisps can be gained in one day. If the Sworn is Dedicated to a Princess of the same court and assists her in a task that fits the court's philosophies then providing he got a single success on the roll he gains a bonus Wisp, this can take him to a maximum of 2 * Dedication + 1 Wisps per day.

A Sworn may buy this merit multiple times for Dedication to multiple Princesses. The daily limit applies to each Princess independently, however only one roll may be made per scene or hour. 

\subsection*{Dream Travel (\textbullet\nobreak\hspace{.16667em}or\nobreak\hspace{.16667em}\textbullet\textbullet\textbullet)}

Prerequisite: Not supernatural\\

\noindent Your character has psychic powers, specifically he can project his mind into the Dreamlands. For the most part this uses the same rules as Princesses, spending Willpower instead of Wisps where appropriate. See the Dreamlands appendix for more details. 

There may be some small differences, some psychics might prefer to use meditation or concentration instead of REM sleep and so gain different situational bonuses to reaching the Dreamlands. The only significant difference is that psychics are not targeted by The Trap and so must discover their abilities unaided. A strong desire to escape regular life or a coincidence where someone with the gift decides to train themselves in lucid dreaming are the most common method. 

The three dot merit allows the psychic to create a shield against the gales using the power of their mind, this provides the same protection as Regalia. The shield can take any form, from nothing at all to transforming into an astral avatar. However it looks the psychic may spend a point of Willpower to manifest their shield for the duration of a trip to the Dreamlands. 

Dream traveling psychics may buy the White Rabbit Dreams merit. It is also not unknown for them to have additional psychic gifts, such gifts usually work in both Earth and the Dreamlands but at the Storyteller's discretion an exp discount may be applied if gifts that typically work in both worlds are modified to only in the Dreamlands. \\

\textbf{Drawback:} Doors allow travel both ways, your character is especially \textit{accessible} to astral threats like Amanojaku. Hopefully he has the training or gifts to defend himself.

\subsection*{Secret Keeper (\textbullet\textbullet\nobreak\hspace{.16667em}or\nobreak\hspace{.16667em}\textbullet\textbullet\textbullet)}

Prerequisite: Player character\\

\noindent You know of a Princess' secret identity, and you've sworn above all else to honor her trust in you. When rolling to protect against any attempt to pry the secret from you, be it mundane or supernatural, you apply the Rote Action rule. The three dot version of this Merit applies to an entire Clique or Nakama. 

Players who wish for their character to have a NPC secret keeper should invest in the True Friend Merit. 

\begin{quote}
HERE MEN FROM THE PLANET EARTH FIRST SET FOOT UPON THE MOON
JULY 1969, A.D.
WE CAME IN PEACE FOR ALL MANKIND
 \end{quote}
 
The classroom of an inspirational teacher. A humble yet gifted artist's studio. The best, hardest working hospital in town and the underfunded free clinic opening against the odds in the most deprived area. A truly welcoming and tolerant church. The laboratory tirelessly working on the latest disease to evolve and the library that preserves wisdom of ages gone by. These are the Blessed places, they stand on a legacy of the very best of humanity.
 
Anyone touched by the Light - any Beacon, Sworn, Shikigami, or Noble - feels her heart lift when she arrives in a Blessed area, though she may not realize why. The Storyteller makes a Perception check, rolling Wits + Composure + Sensitivity - Shadows; if it succeeds the character notices the Light-touched quality of the place she enters.
 
\subsection*{Traits}

A Blessed Place has the following traits:

\subsubsection*{Size}

The Size of a Blessed place is measured on the Sanctuary scale. It’s determined by the scale of the acts that created it - roughly, the number of people directly affected by them. Blessed areas can grow larger, if deeds inspired by those that made it are done close by. The table below defines how large a Blessed Place will grow. 

\subsubsection*{Beauty}

Rated from 1 to 5, and determined by the level of excellence demonstrated by the acts that created the Blessed place. Their excellence may be aesthetic, intellectual or moral. The fame of the acts also influences Beauty; spreading the tale of a great deed far and wide often makes the Blessed place formed by it stronger.

To determine the Blessed Place’s size, find the area or amount of people affected in the table below. To determine the Blessed Place’s beauty, find the class of deed in the table below.

	\bigsidebarstart

 \centering 
    \begin{tabulary}{\textwidth}{|l|L|l|L|}
    \hline
    Size of Blessed Place created & People or Area Affected & Beauty & Class of deed \\ \hline
    0    & One person & 0 & Minor: Creating beautiful art, original and useful research, successful humanitarian works. \\ \hline
    Sanctuary \textbullet    & Neighborhood & 1 & Significant: Creating a masterpiece, a discovery that shapes a discipline, lifesaving and difficult humanitarian works \\ \hline
    Sanctuary \textbullet\textbullet    & Small town & +1 &  Multiple deeds accomplished over a long period.\\ \hline
    Sanctuary \textbullet\textbullet\textbullet    & City & +1 & Deeds required overcoming personal demons.\\ \hline
    Sanctuary \textbullet\textbullet\textbullet\textbullet    & State, small country     & +1 & Deeds are accomplished at a great risk of social censure. \\ \hline
    Sanctuary \textbullet\textbullet\textbullet\textbullet\textbullet    & Continent or world        & +1 & Deeds are accomplished at great risk of injury or death. \\ \hline
    \end{tabulary}
    \bigsidebarend


\subsection*{Aspirations}

A Blessed place has from one to three Aspirations, the types of deeds it inspires people to imitate. These should be generally phrased, and reflect the people responsible for the place’s creation - an artist motivated by his love of art? A scientist dedicated to discovery? A good Samaritan? It’s unusual for a Blessed place to have more than one Aspiration, however great its Beauty; when it happens, it’s generally thanks to a complicated history in which several people were doing very different things, which were nonetheless connected.

While they remain within a Blessed area, any mundane characters treat its Aspirations as their own - they gain a Beat if they fulfill one, their presence may add or open Doors in social maneuvers, and so on. Light-touched characters can do the same, Princesses and Beacons may use the Aspirations as Dreams if they are fitting, while Dark creatures can’t benefit at all; other supernatural beings may be swayed at the Storyteller’s discretion. Since the Aspirations belong to the place they don’t disappear when fulfilled - the place never loses its power to inspire - though each character may fulfill each Aspiration only once a session.

Furthermore, any action taken within the bounds of a Blessed place that directly opposes the fulfilling of any of its Aspirations takes a penalty equal to the place’s Beauty rating. Conversely, any action that directly furthers one of the place’s Aspirations takes a bonus equal to its Beauty. This effect applies to all characters: mundane, Dark, Light-touched, or other. Students feel inspired in the great artist's workshop, while bratty children feel guilty and hesitate in their boasts of one-upmanship within in the tomb of past heroes.  

\subsection*{Corrupting a Blessed Place}

Like so many beautiful things, Blessed Places can be fragile, and need to be looked after. Whenever an act generates Taint within a Blessed Place, the Blessed Place gains Shadows equal to the Severity of the act. If the evil that has been done can be mended in some way (the murderer found and brought to justice, for instance) before the end of the current Story, then the Shadows dissipate; however, each Shadow left on a Blessed Place when a Story ends, reduces its Beauty by 1. If a Blessed Place drops to 0 Beauty in this way, it is immediately destroyed. Any excess Shadows are lost, the last remnants of Light at least preventing a Tainted Place from forming in its stead right away.

\section*{Invocations}
\phantomsection
\label{sec:Invocations}

Invocations are Beliefs of the Queens. A Queen's light shines out onto the world and Princesses reflect that light to shine ever brighter. A Princess may buy up to 5 dots in an Invocation. A Princess may \textit{apply an Invocation} to boost her Charms, normally this means adding her dots in the Invocation to her dice pool for the Charm but some Charms list different effects. In most circumstances applying an Invocation costs a Wisp, but each Invocation has a set of situations which removes the cost. There is also, for each Invocation, a prohibition that a Princess should not break; if she does she loses the right to apply that Invocation for a period of time, she may still use Charms that depend upon that Invocation as a Prerequisite.

Charms that don't require a roll to activate gain power from Invocations in other ways; for example, Barrier Jacket uses an Invocation to increase a Princess' armor. It is never possible to apply two Invocations to a single dice pool, or to stack Invocations in any other way. Some Charms, called Invoked Charms, are linked to one specific invocation. These are denoted by having dots in an Invocation, as well dots in a Charm family. A Princess may not learn an Invoked Charm unless she has sufficient dots in the relevant Invocation, and when using an Invoked Charm only the linked 
Invocation may be used to increase the Charm's potency. 

Applying an Invocation is a reflexive action, and so a Princess may apply multiple Invocations within a single turn if she has the opportunity and has the Wisps needed to do so. She may apply the same Invocation as many times as she wishes.

\subsection*{Invoking For Willpower}

Aside from adding power to Charms, an Invocation can inspire Nobles who strive to live up to its principles. What could be more encouraging than knowing a person of almost divine wisdom supports of your actions? If the Storyteller agrees that a Princess’ actions during a scene reflect the ideals of an Invocation the Princess has learned, and she hasn’t broken its ban, the Princess rolls her dots in that Invocation as a dice pool. (This pool cannot be modified in any way - she may not spend Willpower on it, and no magic or Condition affects it.) She may do this once per session, for each Invocation separately. She need not be transformed to use her Invocations this way.
 
 \begin{myindentpar}{10pt}
    \emph{Dramatic Failure:}  The Princess’ recent actions feel a bit hollow; she doubts the ideals that inspired them. Until the scene ends, she may not spend Willpower on actions that reflect the Invocation.

    \emph{Failure:} The Princess goes on with her life.

    \emph{Success:} The Invocation confirms the Princess’ actions. She regains one spent Willpower point.

    \emph{Exceptional Success:} The Invocation revives the Princess’ flagging will; she regains all spent Willpower.
\end{myindentpar}

\subsection*{Choosing Your Invocations}

The choice if Invocations is an important one, an Invocations shapes empowers a Princess' magic. The easiest trap for a player to fall into is jump to thinking which Charms she wants. Invocations are not a neatly packaged source of power for any Princess' who likes the look of a certain Charm. Rather, an Invocation is the ideals and Beliefs of a Queen given form, resonating in the souls of mankind through the inherent magical nature of a Queen.

When choosing Invocations a player should think about what views and beliefs a Princess shares with a Queen. The more a Princess has in common with a Queen, the more she upholds a Queen's ideals, the brighter the Invocation shines within her. With practice (represented by spending experience) she can draw on that Invocation to enhance her magic. The Queen's philosophies can provide a yard stick. A Princess who only follows one of the Queen's philosophies is unlikely to advance beyond one or two dots in an Invocation. Agreeing with two philosophies can reach three or perhaps even four dots. A Princess who broadly agrees with and lives up to all three philosophies can with time master an invocation. (For Lacrima, a Princess is considered to uphold the philosophies if she applies similar ideals to her own friends, family or some other organization instead of Alhambra.)

And what of the Twilight Invocations? Not every philosophical ideal of the Twilight Queen's is pure evil to be rejected by any right minded Princess. Protecting those you love or opposing the Outer Dark in all it's forms are ideas commonly found among the Radiant. So what stops a Princess waking up one morning and realizing she now has a Dot of Lacrima? The answer is choice. A Princess who believes in protecting those she loves can feel a connection to Lacrima deep within herself and knows she can draw power from it, but she can choose not to. The connection to an Invocation does not shape her beliefs, it is shaped by her beliefs. An Invocation she chooses not to use, or vows to stop using, cannot turn her to the dark side. 






\bigsidebarstart



\section*{Dancing in Twilight}

One of the dirtier secrets of the Radiant is the sheer number of Nobles who, at one point in their life or another, dabble in the Invocations of the Twilight Queens. Dabbling in this case can be anything up to buying two or three dots in the Invocations; more requires a strength of commitment to the ideals of the Queen which is largely incompatible with remaining a member of the Radiant (for one, one's own Queen will tend to object). The truth is that, the flaws in the Twilight Queens are ones which are within the hearts of man as a whole, and so most of the Enlightened have felt the same urges at some point in their life. Note that, however, such dabblers still consider themselves part of the Radiant, and, in the case of those who feel affinity for the Queen of Mirrors, often consider the The Child Queen to be Radiant. They still oppose the Darkness, and will fight against the loyalists of Tears and Storms. One of the greatest topic of debate in the multiple, disparate groups of the Enlightened is how to 
respond to people who use Twilight Invocations, and whether or not they're doing anything wrong.



The Invocation of Lacrima, tool of the Last Empress, the Queen of Tears, is an object of fear and mistrust to the vast majority of the Radiant. It is apparently evil, in a way that neither Specchio nor Tempesta are, because it's Charms are an inherent Compromise of one's Belief, and to many, it is irrevocably tainted by the fact that it has the Charms used to drain Light from the world or transform oneself into a creature of the Dark. But, nevertheless, there are those Princesses who learn it, and who do not follow the Queen of Tears. One of the fundamental things about Lacrima is that it can always be used to protect those you care about, and to a Princess who wants to do exactly that, the expediency of methods means that Lacrima can look attractive. When a friend died because you pushed your views on someone and so Legno was sealed off, the universal applicability of Lacrima for protecting your own starts to look very tempting. The dabblers are typically more experienced, and more able to justify to 
themselves the 
moral compromises involved, in part because although there are Lacrima Charms which are Belief Compromises, dabblers tend to only have access to the ones which are Compromises against very high Belief or not Compromises at all, and so there are less of an objection. Moreover, there exists a notable sub-faction of the followers of the Queen of Diamonds who specifically learn Lacrima for its use in studying the dead and the Underworld, and, aggravatingly to some others, the Queen has even given her support to the practice. Such dabblers often have a strict moral code restricting their use of the Invocations to only using it to study the dead - never to harm, and they point to the gains in knowledge that they have provided to the Radiant; their detractors point out that once the power is there, it can be hard to resist. 

The Invocation of Tempesta, tool of the Seraphic General, the Queen of Storms is for most of those who dabble in it, purely a weapon. It is something that is learned as a tool against the Darkness, and it is used for that role. It does not (generally speaking) heal or right wrongs, but it does eliminate the followers of the Darkness and of Alhambra. In Atlanta, in the mid-90s, almost an entire generation of Radiant Princesses knew at least the basics of Tempesta. So hard-pressed were they by the Darkness and a major, well-organized Alhambran Outpost that it became unofficial policy for elder Enlightened to guide each newcomer towards it, because the universal applicability and raw power of the Invocation in such circumstances was too much of a boon. That large numbers of Nobles in the city were driven towards extremism was viewed as an acceptable price, even if by the end, and the arson of the Alhambran Outpost, there were Enlightened just as callous as the worst of Storms, still officially among the Radiant.
 For those Princesses in the know, too, Tempesta can also be seen as the most moral of the Twilight Invocations, because the innate costs to it, the damage it does to the user, is known, and only affects the self. The common consensus is that it is dangerous, and a sign that a Noble is starting to spiral into extremism, but whether the use of it should be enough to ostracize one of the Radiant is a lot more contentious, especially since in areas where it is done, there are frequently not the Enlightened to spare.

The Invocation of Specchio, tool of the Crown's Custodian, the Queen of Mirrors, however, is the most subtle and insidious of the Twilight Queens, and the influence it is. In part, it is because of the lack of an overall enemy that the Invocation is linked to; while both the Ravens and the Furies provide a highly visible \textquotedblleft there but for the grace of the Light I go" example which other dabblers can keep away from, the egotism and self-righteousness implicitly promoted by the Queen of Mirrors, and the lack of a solid adversary is hard to fight against, especially since many Princesses have tendencies in that direction naturally. Dabblers can slip into following the Queen of Mirrors much more easily, especially since the Queen does not demand fidelity, which means that some elder Radiant are much more harsh in opposing its use. There are those who argue that, unlike Tempesta and Lacrima, use of this Invocation makes you a de facto follower of Mirror's agenda. And on the other hand, there are 
those who count the Queen of Mirrors as among the Radiant Queens, despite the objections of the other five. Moreover, there us \textit{also} a faction in the Court of Diamonds who study Specchio for the unique sources of information it provides, or even to study the strange achronal or spacewarping Charms unique to Specchio, once again the Queen has given her consent. Even outside the Court of Diamonds there are Radiant who find these Charms to be exceptionally useful and irreplaceable with more moral Invocations, not just an improvement over them. In the Court of Diamonds these dabblers protect themselves from the effects of madness by only using Specchio while supervised by a Princess who knows the Reclaim Charm. Of all the Twilight Invocations Specchio might be the easiest to learn, too; all one needs for it is self-belief in ones own cause, and a spark of the brilliance needed to reshape the world. And neither are rare among the Enlightened. 

\bigsidebarend



\section*{Charms}
\phantomsection
\label{sec:Charms}
The basic magics of the Hopeful are called Charms. Each is a unique ability available to the Nobility, some are obviously supernatural, some are subtle, but all are truely magical. 

The Charms are divided into 10 trees. Each Calling has an affinity to 3 of them, learning the Charms within them more easily. 

Each Charm has an Action attribute, giving the type of action using it is:

\begin{itemize}
\item Instant: Activating the Charm is an instant action, taking a turn. The dice pool to roll for the action is noted.
\item Full turn: Activating the Charm is an instant action, but it also requires forgoing Defense in the turn you activate it.
\item Reflexive: Activating the Charm is a reflexive action; it doesn't interfere with your normal action and can be done at any point in the initiative cue. A dice pool to roll is noted.
\item Reactive: Activating the Charm can be done at any point in the initiative cue but it consumes your action for the turn. A dice pool to roll is noted.
\item Permanent: The Charm enhances another ability. No action is required to activate it -- if the proper conditions obtain, it just works. A Princess can always disable a permanent Charm if she so chooses. If no other method is presented a Transformation action will activate or deactivate any number of permanent Charms. 
\item (one of the above) and resisted: The Charm has a target, and you must subtract one of his traits from the dice pool.
\item (one of the above) and contested: The Charm has a target, who reflexively rolls an opposed dice pool; you must roll more successes than he does for the Charm to take effect.
\end{itemize}
    
Other than the permanent Charms, each Charm has a duration, ending when the stated time runs out; and a cost, usually in Wisps. Powerful Charms call for Willpower points; many Charms favored by the Queen of Storms inflict damage on their user, which resists magical healing; and using some Charms favored by the Queen of Tears transgresses against Belief. 

A Charm with a duration of \textquotedblleft lasting” changes its target permanently in a flash of magic but the Charm itself ends immediately. Charms with an \textquotedblleft indefinite” duration leave a continuing trace of magic on their target (often a Condition or Tilt) which other supernatural powers can detect or influence. A Charm with a duration of \textquotedblleft concentration” needs its user’s active attention, and ends if the user spends her turns action on anything but maintaining the Charm.

If a Charm is used on a person, a Princess can use it on herself Charm specifically states she cannot. 


\subsection*{Learning Charms}

Each Charm has a rating in dots, which sets the price and difficulty of learning its basic effect. While Charms need not be learned in any fixed sequence, the more advanced Charms can't be learned without practice with simpler powers. To learn a Charm at its basic level, a Princess must have a sufficient number of dots in other Charms in the same tree (counting each Upgrade as an additional dot), as defined in this table. 

\begin{mybox}

    \begin{tabular}{|l|l|}
        \hline
        Charm Rating & Dots Required \\ \hline
        1            & 0             \\ \hline
        2            & 1             \\ \hline
        3            & 3             \\ \hline
        4            & 6             \\ \hline
        5            & 10            \\
        \hline
    \end{tabular}

\end{mybox}
    
\textit{Example: Before she can learn a Govern Charm rated at three dots, a Princess would have to know: Three one dot Govern Charms, a one dot and a two dot Govern Charm, two one dot Govern Charms one of which has an Upgrade, or a single Govern charm which has two Upgrades.}

A Princess may also treat Invoked Charms (see below) as a pseudo tree for the purposes of learning Charms. For example: If a Princess knows three one dot Perfect Charms that each require Acqua, she is able to acquire any one, two or three dot Charm that has an Acqua prerequisite. She may not acquire an unInvoked Charm that has an Upgrade which depends upon Acqua, however if she buys an Upgrade to an unInvoked Charm that requires Acqua that counts as another dot towards learning further Acqua Charms. 

\textit{Example: Before she can learn Pierian Spring, a two dot Learn Charm that is Invoked by Acqua, a Princess would have to know any one dot Learn Charm, or any Charm or Upgrade that has an Acqua prerequisite.}

A few Charms have another Charm as a prerequisite; to learn these Charms, a Princess must have learned the prerequisite Charm, as well as meeting the total dots required for its basic effect. Finally, your Inner Light limits the dots you can put in a Charm; a Princess cannot learn a Charm with more than Inner Light + 2 dots. 

\subsubsection*{Upgrades}

Many Charms have Upgrades which enhance the Charms effects or expand the Charm into new areas. Upgrades may have prerequisites and each Upgrade increases the effective rating of a Charm by one for the purposes of a Princess' maximum Charm rating as determined by her Inner Light. So a Princess with one dot of Inner Light could not apply an Upgrade to a three dot Charm and could only apply one Upgrade to a two dot Charm. However she could learn multiple upgrades for a two dot Charm; she would simply be limited to applying one Upgrade at a time. 

Many upgrades change the cost or duration of their Charm, add a modifier to the activation roll, or have a Drawback, or trigger Specchio's drawback. Therefore, you are not required to apply any upgrade when using a Charm. When Charms require an activation roll the Princess can choose which Upgrades to use at the point she invokes the Charm. While for Charms that are always on she can enable or disable upgrades as she desires with a Transformation action.

\subsection*{Invoked and General Charms}

By default a Princess can learn any Charm, providing she knows enough dots in the relevant Charm family. Once she knows a Charm she may apply whichever Invocation she wishes to increase the Charm's potency. 

A fair number of Charms work differently. These Charms, known as Invoked Charms, require dots in an Invocation to purchase. They can be recognized by the presence of Invocation dots along side the Charm's dots. Invoked Charms cannot be learned until the Princess has the listed dots in the linked Invocation and cannot be cast with any Invocation but their prerequisite. When using such a Charm, a Princess can apply the prerequisite Invocation to gain its bonus (if conditions permit) or she can make the activation roll without any bonus, but she may not apply any other Invocation's bonus. 

A Princess who breaks an Invocation's ban may still use that Invocation's Invoked Charms. Naturally she cannot apply any Invocation for bonus dice. 

There are also general Charms with upgrades that require Invocation dots -- applying such an upgrade converts the Charm to an invoked one -- and some invoked Charms with upgrades that require more Invocation dots than the basic Charm does. Upgrades that require different Invocations are not compatible; you may not apply both of them to the same activation of a Charm. 

\subsection*{Charm Modifiers}

\subsubsection*{Intimacy}

Some Charms can be used on distant targets by taking penalties to their activation roll.  If a Charm's dice pool is modified by Intimacy, the strength of the emotional connection between the Princess and her target affects the power of the Charm. Consult this table for the penalty to the activation roll: 

	\begin{mybox}

 \centering 
    \begin{tabulary}{1\textwidth}{|l|L|}
        \hline
        Penalty & Strength of Connection                                                                            \\ \hline
        0       & Sensory: You can see the target directly.                                                         \\ \hline
        -2      & Love or hate: The target is a longtime friend or enemy, a member of your family, a romantic lover, hurt you badly, or a prized possession.       \\ \hline
        -4      & Friendship or enmity: You know the target well; or the target is something you own and care about.                         \\ \hline
        -6      & Compassion or dislike: You know the target slightly; or the target is an a typical item you own.         \\ \hline
        -8      & Acquaintance: You have met the target briefly; or the target is an item you have held.  \\ \hline
        -10     & Few feelings: You have a verbal description of the target, but nothing more.                         \\  \hline
		-       & Unknown: You cannot affect the target via Intimacy. \\
        \hline
        \end{tabulary}
    \end{mybox}


A Princess does not have to use her own emotional connection. By touching another person she may use their Intimacy instead of her own. This increases the Intimacy modifier by one step (add the intermediary's Composure as an additional penalty if they are unwilling). A Princess can also use the connection between a person and a treasured object, this increases the penalty by two steps. 

If the Princess has a keepsake, a treasured object that symbolizes her relationship with another (such as a love letter) this reduces the penalty by one die. If she uses somebody else's connection she can reduce the penalty by one die if she touches the that person and one of their keepsakes. The benefit does not stack with multiple keepsakes. 

If a character is targeted via Intimacy and they are aware of it they may respond by using any of their own abilities that work via Intimacy at the same dice pool modifier. If a Princess witnesses someone or something else being targeted via Intimacy they may respond using the original dice pool modifier plus one or two steps as is normal for using a proxy. In essence, the character sends their magic back down the channel someone else created through Intimacy. At the Storyteller's discretion this can work with mechanics that work similarly to Intimacy, such as a Mage's Sympathy. 

\subsubsection*{Sanctuary}

Several Charms affect large areas, and the effort required to invoke most of those Charms scales with the affected area. These charms will refer to the target area using the keyword Sanctuary. 

Charms that use the Sanctuary keyword are typically easier to activate when the area in question is one the Princess values, and conversely are difficult to use in areas she dislikes or fears. Consult the following suggested modifiers: 

\begin{itemize}
\item +5 Exceptionally strong memories of love and safety, such as a childhood home
\item +3 Strong memories of love and safety, such as a long time friend’s home
\item +1 The Princess owns the property
\item +1-3 The area is a Safe Place for the Princess (bonus = half the Merit dots, rounding up)
\item -2 General feeling of dislike, such as a school room where the Princess was board.
\item -3 The Princess has never been in the area before
\item -5 Strong memories of fear or horror
\end{itemize}


\begin{mybox}

 \centering 
    \begin{tabulary}{1\textwidth}{|l|L|}
        \hline
        Rating  & Sanctuary Size \\ \hline
        X       & Just enough room for one person to stand or lie down (Typically Charms using the Sanctuary keyword can’t affect an area this small) \\ \hline
        \textbullet      & A small apartment or a cabin in the woods; 1-2 rooms   \\ \hline
        \textbullet\textbullet      & A  large apartment or small remote house; 3-4 rooms   \\ \hline
        \textbullet\textbullet\textbullet     & A converted church, warehouse or large house; 5-8 rooms  \\ \hline
		\textbullet\textbullet\textbullet\textbullet  &    A mansion or apartment building; 9-15 rooms  \\ \hline
		\textbullet\textbullet\textbullet\textbullet\textbullet    & A vast palatial estate, a city block, a skyscraper or a small village \\ \hline

    \end{tabulary}

\end{mybox}


\subsubsection*{Commonality}

Some Charms are modified by Commonality and can be applied to an entire social group at no additional cost. Commonality does not increase the range of a Charm and cannot be combined with Intimacy, so only the members of the targeted social group who are present will be affected. If such a Charm is resisted, the member of the group with the best resistance trait resists for the whole group; for a contested Charm, each member rolls to contest individually but the Storyteller may combine rolls for groups of NPCs to keep the game moving along. The difficulty of extending the Charm over a social group depends on the commitment the members have towards each other, or towards the common goal for which the group exists.  Consult this table for the appropriate penalty: 

\bigsidebarstart

 \centering 
 \begin{tabulary}{1\textwidth}{|l|L|}
        \hline
        Penalty & Strength of Connection                                                                                                                                                                        \\ \hline
        0       & Individual: one person, the default.                                                                                                                                                          \\ \hline
        -3      & Dedicated: everyone in the group affected is mutually friends, or the group's goal takes up the majority of its members' time and thought; an established nakama, a family.                  \\ \hline
        -6      & Concerned: everyone in the group knows each other, or all the members make significant sacrifices for the common goal; a new nakama, the long term employees of a small business.  \\ \hline
        -9      & Interested: everyone in the group is mutually acquainted or better, or the members share an interest in the group's goal; a social club, employees of a large corporation.                    \\ \hline
        -12     & Casual: some members of the group have only encountered each other; the group is merely a lot of people gathered in one place, with no common interest.                                       \\
        \hline
        \end{tabulary}
\bigsidebarend

\subsection*{Extended Actions and Charms}

Charms do not exist in isolation, they are deeply tied into a Princess' hopes and beliefs. Usually a Princess' Charms are influenced by her hopes, but a particularly strenuous act of magic can have a noticeable if fleeting affect upon a Princess' emotions. 

If a Princess performs a Charm or uses an Embassy Privilege as an Extended Action she may gain the following conditions, in addition to any effects of the Charm itself: 

\begin{itemize}
 \item Exceptional Success: Luminous
 \item Failure: Stuck Magic
 \item Dramatic Failure: Shaken
\end{itemize}

\subsection*{Teamwork Actions and Charms}

Princesses are very good at working together and all Charms with an activation roll can be invoked as a Teamwork action, just so long as every participant knows the base Charm. Only the primary actor needs to spend Wisps, Willpower, or take Damage in the case of Tempesta Charms; however when invoking a Lacrima or Specchio Charm (including an Invoked Upgrade) the Invocation's Drawback applies to all participants. All participants in a teamwork action must apply the same Invocation or no Invocation at all, and secondary actors must still pay a Wisp for the Invocation unless they qualify to apply the Invocation without a cost. 

The Merit Troupe Magic removes the need for secondary actors to know the base Charm, but only for Inspire Charms. 

\subsection*{Clash of Wills}

When two Princesses (or anyone with magical powers) attempt to invoke conflicting magics a Clash of Wills occurs. This could be two Hope Engineers fighting to control the same body of water, a Wild trying to magically calm a crowd while a Fury tries to whip that crowd into a murderous fury or a group of Jewels each trying to inspire people to believe in their political party. 

When a Clash of Wills occurs each party makes a contesting roll to resolve the clash, the winning party's magic takes precedence. The dice pools are as follows: 
\begin{itemize} 
\item Princess (including Ghost Princesses): Inner Light + half Belief (rounding up)
\item Sworn or Beacon (with a Bequest): half Integrity (rounding up)
\item Shikigami: half Integrity (rounding up)
\item Dreamlanders (including Wardens but not the former inhabitants of the Kingdom): Power + Rank
\item Goalenu: Power + Rank (unvesseled) or Attribute + Rank (vessled)
\item Creature of the Darkness: Shadows
\item Dethroned: Inner Darkness
\end{itemize}

Any character may spend a point of Willpower to bolster their Clash of Wills roll. Charms and other supernatural abilities with long durations are more durable and more likely to succeed in a Clash of Wills. If the effect lasts a day add one die. For a week long effect add two dice, three for a month, and add for dice if the effect lasts a week or longer. 

\subsection*{Disabling Charms}

Some Charms have Upgrades that provide passive or automatic abilities. When a Princess Transforms the player may choose which Charms are in effect and a further Transformation action may be made at any time to switch which Charms are in play. This also applies to individual Upgrades. 

\subsection*{Regalia}

Some permanent Charms add equipment to the Princess' Transformed self, or give her the ability to add equipment. This gear forms the Princess' Regalia. When a Charm states that an item is part of the Regalia it can appear on her person, ready for use, when the Princess transforms if she wishes it to. A single piece of Regalia can hold the power of multiple charms if the player wishes. If a Charm does not state that it creates Regalia a player may still declare that it does so, but this has no mechanical effect. 

A Princess can switch any number of Regalia pieces into or out of solid existence with a Transformation action; the dice pool for this action is the same as the one for full Transformation. However, if she returns to mundane form she must dismiss all her Regalia as well. Naturally, pieces of Regalia help a Princess with her goals only when they are in existence - any bonuses they confer are lost if she sends them away.

All Regalia has a Structure equal to Belief and a Durability of Inner Light. If a piece of Regalia is destroyed or lost, the Princess can recreate or summon it while transformed; as with a dissolved phylactery, she must spend a Wisp and a Willpower point to recreate a piece of Regalia, which takes a full round and this requires her absolute concentration. 

If another was to steal a Princess' regalia they would be unable or have great difficulty in using it, but they would be able to target the Princess through an unbreakable Intimacy link with a strength equivalent to love. Fortunately any technique that would cut the connection Princesses use to summon Regalia would, by definition, destroy the Regalia.  


\subsection*{Appear}
\phantomsection
\label{subsec:Appear}

The Appear Charms change your appearance, and produce other illusions. Seekers and Troubadours have affinity for them. 

\subsubsection*{Flair Emotion (Appear \textbullet)}

\noindent \textbf{Action:} Instant 

\noindent \textbf{Cost:} 1 Wisp 

\noindent \textbf{Dice Pool:} Unrolled 

\noindent \textbf{Duration:} Lasting \\

\noindent The Princess takes a Wisp's worth of emotional energy and releases it into the air. To psychic empaths, Princesses using Passion’s Light, and similar abilities it shines like a flair and can be used as a signal. Usually no roll is required, the Princess spends the Wisp and barring a large distraction any empath in the area will notice, but a Composure + Empathy roll may be used to get into the right emotional state if the Nakama have agreed to use specific emotions as specific signals. 

An alternative use for this Charm is to release the Wisp in close proximity to an empath, this hits like a flash-bang and triggers a Presence + Empathy vs Composure + Supernatural advantage roll. If the Princess wins the roll she inflicts an appropriate Tilt for either a turn or one turn per success. The Tilt chosen depends on how the target senses emotions and the duration depends on the severity of the Tilt. A character who hears emotions might take Stunned for one turn, or Deafened for several turns. 


\subsubsection*{Grey Masquerade (Appear \textbullet)}

\noindent \textbf{Action:} Instant

\noindent \textbf{Dice Pool:} Manipulation + Empathy

\noindent \textbf{Cost:} 1 Wisp 

\noindent \textbf{Duration:} 24 hours \\

\noindent Princesses enjoy powerful protection to keep their secret identity secret but not everyone is so lucky. Beacons cannot hide the unusual amount of Light they produce by storing it inside a Transformed self. Sworn are visibly under an enchantment. So helpful Princesses can use this Charm to disguise the signs. After all, if your goal is to give your sister the magic she needs to fly away from danger painting a target on her back defeats the purpose. 

This Charm can only be used on Beacons, Sworn, or Shikigami. It is commonly known by Shikigami who travel independently of their Noble sponsor to seek out Princesses to mentor. 

\begin{myindentpar}{10pt}\emph{Dramatic Failure:} The Princess accidentally sends out a flair of magic, perceivable over a reasonable distance by anyone who can sense such things. 

    \emph{Failure:} The target's supernatural nature raiment.

    \emph{Success:} The target takes the Masquerade Condition, any passive supernatural power which would detect their supernatural nature, such as the Charm Reflected Light, will automatically fail. Active supernatural powers, such as the Charm A Magical Girl's Eyes, must overcome a Clash of Wills against whoever invoked this Charm. If it cannot overcome the illusion the target registers as a completely normal human (or a completely normal something else in the case of Shikigami). This Charm provides no protection to active uses of power by the Charm's subject. If a Sworn activates his Practical Magic near someone with Reflected Light they will feel it as normal.

    \emph{Exceptional Success:} The duration lasts a whole week.
	
\textbf{Upgrade: Masqued Ball}

\textbf{Cost:} +1 Willpower

The Princess can mask the enchantments she or another Princess created, hiding it from discovery through abilities like A Magical Girl's Eyes. She could invoke Masqued Ball upon artwork enchanted by Living Image, a character under the effects of a Bless Charm, an inanimate object with the Held Charm Condition, or any other magic that lingers after the invocation of a Charm, Privilege, or Edict. For the duration of the Charm all Noble magic affecting that object or person is hidden. As with regular uses of Grey Masquerade the illusion can be penetrated through a Clash of Wills. 

If the Charm the Princess wishes to hide works on a Sanctuary scale the roll becomes an Extended Action with a target number equal to twice the Sanctuary's size. If the Princess hides the Consecrated Condition then the duration extends until the end of the Consecrated Condition and it will hide all Charms that are built atop the Condition. It will not hide other Charms that happen to be within the bounds of the Consecrated area, such as a person under the effects of a Bless Charm relaxing inside a Consecrated house. 

\textbf{Upgrade: Close the Third Eye (Acqua \textbullet\textbullet)}

By studying esoteric lore of the human mind the Princess can learn to shield a mortal psychic's signature, hiding Supernatural Merits such as Psychokinesis from magic that might sniff out a psychic. This changes the dice pool to Wits + Occult, and allows the Princess to invoke this charm on a pure Mortal with an appropriate Merit. Apart from that the Charm works as normal.  

\textbf{Upgrade: Gold Masque (Aria \textbullet\textbullet)}

\textbf{Cost:} +2 Wisps

The Princess' ability to mask a supernatural nature grows more potent, strong enough to hide a Transformed Princess. 

\textbf{Upgrade: Perennial Veils (Lacrima \textbullet\textbullet\textbullet)}

\textbf{Cost:} +1 Willpower Dot

The Charm's Duration extends to Indefinite.

\textbf{Upgrade: Funerary Masque (Lacrima \textbullet\textbullet\textbullet\textbullet)}

The Princess' twists her ability to conceal the supernatural, and instead conceals the mundane from supernatural eyes. Any supernatural being must make a Clash of Wills in order to notice the target. If they fail the contest they will only remember the target as a generic unimportant person and will never directly interact with them. A hungry Darkspawn who's willing to eat any human it can find will fail to consider the target even if no other options are available. Princesses are a special case to this rule, in their mundane form they can see a person under the effects of Funerary Masque as clearly as any mortal, but in their Transformed identity they must make a Clash of Wills to penetrate the veil. 

Funerary Masque can only be invoked upon Mortals and Beacons. 

\end{myindentpar}


\subsubsection*{Light the Way (Appear \textbullet)}

\noindent \textbf{Action:} Instant 

\noindent \textbf{Cost:} 1 Wisp 

\noindent \textbf{Dice Pool:} Unrolled 

\noindent \textbf{Duration:} Inner Light or Invocation hours \\

\noindent The Princess can, at any time, create a soft light that banishes the thickest natural darkness. The light source can take any source the princess wishes, from a torch or a glowing halo to a simple ambient glow. She can bring the light up or douse it reflexively, and set it to shine at any intensity short of full sunlight.

While this Charm is for seeing in dark environments it can be used as an emergency defense. The Noble can also focus the light to a bright, narrow beam into someone’s eyes to inflict temporary blindness. This is a normal attack using Dexterity + Athletics, and the target’s Defense applies. On a normal success
the target is Blinded [CofD 281] in one eye; on an exceptional success he is Blinded in both eyes. Either way, the Tilt lasts for one turn.

\begin{myindentpar}{10pt}\textbf{Upgrade: Dazzling (Fuoco \textbullet)}\\

\noindent The Noble can reflexively make her light burn at a constant eye-watering intensity, blinding anyone nearby who looks at her, by spending a Wisp. Anyone who can see her must succeed on a Stamina + Composure roll, penalized by the Noble’s Fuoco, or be Blinded in one eye until he looks away from her, or she dims her light again, whichever comes first.\end{myindentpar}

\subsubsection*{Twenty Faces (Appear \textbullet)}

\noindent \textbf{Action:} Instant

\noindent \textbf{Dice Pool:} Wits + Subterfuge

\noindent \textbf{Cost:} 1 Wisp 

\noindent \textbf{Duration:} 1 scene \\

\noindent As the purest expression of a Noble soul a Princess' regalia and heraldry can look like anything, but it's almost certain to stand out in some way when compared to outfits made of mere physical matter. Even a Princess of Hearts who wears Savile Row style suits will have some some \textit{je ne sais quoi} about her. Not enough to gain mechanical advantages (unless represented by taking Merits like Striking Looks as single form Merit) or be outed as a supernatural being, but enough to be memorable. And sometimes standing out isn't what the Princess wants. With this charm she can hide her heraldry and adorn herself with appropriate clothing to blend in wherever she goes.

\begin{myindentpar}{10pt}
    \emph{Dramatic Failure:} The Princess cannot use Twenty Faces for the rest of the scene.

    \emph{Failure:} The Princess' appearance and raiment are unchanged.

    \emph{Success:} The Princess wraps herself in an illusion that conceals her heraldry; she seems to be dressed in appropriate clothes for the area. In an office she will manifest a suit or business causal, on a military base she will acquire a uniform. The illusion does not alter her features - anyone who has seen her transformed self will recognize her if they meet her - but her status as a Noble is hidden. Supernatural powers that pierce illusions must beat the Princess' successes to see her true appearance.

    \emph{Exceptional Success:} The illusion blurs the Princess' features. People who have met her transformed before don't recognize her, unless they make a special effort to identify her (succeeding on a Wits + Investigation roll.)
	
\textbf{Upgrade: Fashionable}

The Noble can choose to stand out from the crowd rather than blending in. She may reshape her heraldry into high fashion [CofD 277] suitable for her current venue, or one she’s about to attend. The equipment bonus of her new apparel is half the applied Invocation or half Inner Light (round up) to a maximum of +3. 

If she stays within the general theme of her normal Regalia (colors, decoration, motifs etc.) and is recognizable as herself, the Noble may use Charms and other Regalia without resetting her clothes, but doesn’t grant the benefit of appearing mundane. 
	
\textbf{Upgrade: Disguised}

\noindent \textbf{Cost:} +1 Wisp\\

\noindent The Noble can change her features dramatically. By spending 1 Wisp, she may alter all aspects of her appearance - ethnicity, height, build, even sex - within the limits of her Size trait. (A Princess with the Giant Merit will always be a very tall man or woman, no matter how this Charm is used; a child Princess can look like an adult with dwarfism but cannot significantly increase her height, and an tall Princess cannot pass herself off as a child.) All mundane attempts to see through the illusion fail automatically, and supernatural powers that pierce disguises start a Clash of Wills. Applying Disguised does not give her extra protection from being mundanely detected as an outsider or impersonator. If she disguises herself in a punk outfit then attends a black tie dinner she'll still stick out as a sore thumb. 

The Princess can even attempt to take on the appearance of a specific person, applying the Intimacy modifier for her connection to the person she wishes to impersonate. The Princess and her target must still have the same Size.\\

\noindent \textbf{Drawback:} Any use of Charms dismisses the illusion. Even when disguising herself as another Princess, a Noble's magic is so distinctly hers that it is incompatible with this kind of disguise.  \\

\textbf{Upgrade: Mantle of Maturity}

\noindent \textbf{Prerequisite:} Disguised\\

The Princess can change her apparent size to anywhere between five and six feet, allowing a child Princess to disguise herself as an adult. The Princess' actual size doesn't change, meaning that her real physical body will not match up to her illusory body. People who touch her illusory body will feel an illusory solidity and prevent themselves from passing their hand through the illusion. However if they are moving at speed they might not be able to stop themselves pushing through the illusion and any thrown object will pass straight through the illusion. The effects of this are unique to every Princess: One Princess' illusory body might distort like an out of tune television, another might burst into butterflies that fly around the obstruction before reforming. Regardless of the exact effect, it will always expose the illusion as supernatural. 

\end{myindentpar}

\bigsidebarstart

\subsubsection*{Flaunting Frills at Formal Functions}
 
 Elaborate outfits are one of the most recognizable features of the magical girl genre, so of course Princesses have their heraldry too. But what happens when a Princess walks into a town hall meeting while dressed in frills and lace? 
  
 Given the how central costumes are to maho shojo your group might feel it defeats the point if they have to constantly replace or change out of their iconic outfit. In this case Princesses could have a supernatural trait that makes their Heraldry seem normal. It could be a product of a Noble's own innate magic or an extension of the unique psychological traits that make mortals in the World of Darkness try to keep their head down and ignore the supernatural. The Twenty Faces Charm remains useful for when the players need to disguise themselves as a specific role or change the image they project. The public may ignore the fact that a Princess of Diamonds is dressed like she's stepped out of science fiction, but people will still notice her outfit is rather nerdy. 
  
 However if you prefer a mortals could react to a Princess' appearance as they would react to anyone else who dresses the same way; depending on your preferences this could either require Princesses to carry spare clothes with them or to use Twenty Faces if they wish to avoid awkward questions when interacting with people outside the know. As some completely valid styles of Heraldry, such as the Embassy to the Economy's fashion for suits and ties, become mechanically advantageous under this style of play the Storyteller may wish to impose a Merit requirement on discrete Heraldry to keep things fair. It would go against the themes of Princess if the frilly dresses where mechanically disadvantageous. 
  
 Regardless of your group's choice the Supernatural Heraldry Condition can be taken for a character who's heraldry draws far more attention than the average Princess'.
 
\bigsidebarend

\subsubsection*{Gone With the Wind (Appear \textbullet, Aria \textbullet )}
\noindent \textbf{Action:} Reflexive

\noindent \textbf{Dice Pool:} Dexterity + Stealth vs Wits + Composure

\noindent \textbf{Cost:} 1 Wisp 

\noindent \textbf{Duration:} Lasting \\

\noindent For a few crucial seconds the Princess becomes elusive and all eyes slide off her. This Charm may be invoked reflexively when someone is attempting to track or follow the Princess. Each subsequent attempt to use this Charm on the same person in a scene grants a cumulative +3 bonus to their Contesting roll.

\begin{myindentpar}{10pt}\emph{Dramatic Failure:} The Princess attracts attention. The target may reroll to track the Princess and takes the better result. 

    \emph{Failure:} The Princess' appearance and raiment are unchanged.

    \emph{Success:} The target is suddenly dazed and cannot spot the Princess, he must scramble to find her again. All successes on his last roll are removed but he may try again the next turn, assuming the Princess hasn't used this that time to vanish. If the Princess is attempting to pick the target's pocket, or some similar instant action, this essentially gives her a second try at a Contested roll. If the target beats her mundane roll then so long as she has a single Success she can use this Charm to erase her opponents successes and win the mundane roll. 

    \emph{Exceptional Success:} In addition to the usual effects the Princess' first Stealth roll in the next few turns gains 8again\end{myindentpar}

\subsubsection*{A Minute of Silence (Appear \textbullet, Lacrima \textbullet )}

\noindent \textbf{Action:} Reflexive, unrolled or Instant,

\noindent \textbf{Cost:} 1 Wisp 

\noindent \textbf{Dice Pool:} Dexterity + Stealth vs Wits + Composure

\noindent \textbf{Duration:} 1 scene \\

\noindent The Princess spends a wisp and shrouds herself in silence. All attempts to locate her by listening fail automatically, though the right supernatural power may trigger a Clash of Wills. Moreover, she may spend a Wisp and roll Dexterity + Stealth vs Wits + Composure to make an object or person she touches totally silent for a number of turns to her Lacrima. 

\begin{myindentpar}{10pt}\emph{Dramatic Failure:} A banshee like wail emanates from somewhere close to the Princess for a turn.

    \emph{Failure:} Nothing happens. 

    \emph{Success:} The Princess inflicts the Silenced Tilt on her target. 

    \emph{Exceptional Success:} If targeting a living being the Princess also inflicts the Shaken Tilt through the sheer wrongness of not being able to hear one's own body. \end{myindentpar}

\subsubsection*{Tranquillise (Appear \textbullet, Legno \textbullet )}

\noindent \textbf{Action:} Instant

\noindent \textbf{Cost:} 1 Wisp 

\noindent \textbf{Dice Pool:} Composure + Survival

\noindent \textbf{Duration:} 1 scene \\

\noindent The Princess wraps herself in an illusion of a peaceful glade. 

\begin{myindentpar}{10pt}\emph{Dramatic Failure:} The Princess takes the Sensory Overload Condition and is considered to be in an uncomfortable environment until the scene ends. 

    \emph{Failure:} Nothing happens. 

    \emph{Success:} For the remainder of the scene, or until she ends the Charm, the Princess can hear nothing but a gentle stream in the distance, she smells nothing but a spring breeze, and the temperature feels cool and dry. If the environment is extreme enough to cause actual harm the Princess will still feel pain even though she feels like she is at a comfortable temperature. This Charm is usually used to remove situational penalties when trying to study, meditate, or fall asleep.

    \emph{Exceptional Success:} For the duration of the Charm the Princess takes the Blessing Condition and adds +2 to all meditation rolls. \end{myindentpar}

\noindent \textbf{Drawback:} The Princess takes the Deafened tilt for the duration of the Charm, and cannot smell either. 

\subsubsection*{Zephyr’s Track  (Appear \textbullet, Aria \textbullet )}

\noindent \textbf{Action:} Instant

\noindent \textbf{Dice Pool:} Dexterity + Stealth \\

\noindent \textbf{Cost:} 1 Wisp\\

\noindent \textbf{Duration:} Scene/Lasting \\

\noindent The Noble floats through the world like a breeze, leaving no trace of her passing except by her own will.

\begin{myindentpar}{10pt}

\emph{Dramatic Failure:} The Noble becomes clumsy, prone to slip and make noise when she wants to hide. She takes a -2 penalty on all Stealth dice pools until the scene ends.

\emph{Failure:} The Noble leaves traces of herself behind her as usual.

\emph{Success:} Gentle puffs of air follow the Noble, cleaning away any traces the Noble leaves in the wake of her movements. Footprints and fingerprints are brushed away; hairs and skin cells float behind her without settling; drops of sweat or blood evaporate. Any attempt to follow her trail or deduce her actions from physical evidence fails. This wind follows the Princess for one scene, but it's effects are Lasting. 

\emph{Exceptional Success:} The Noble can choose to leave false evidence behind her that suggests a person quite unlike herself passed where she has gone.

\textbf{Upgrade: Multiple}

Modified by Commonalty

The breezes can clear away traces of many people at once. The Noble uses the Charm on the members of an organization she can see, using the Commonalty modifier. If she succeeds, all traces of their passing are erased.

\end{myindentpar}

\subsubsection*{Phantom (Appear \textbullet\textbullet)}
\noindent \textbf{Action:} Instant

\noindent \textbf{Dice Pool:} Wits + Expression \\

\noindent \textbf{Cost:} 1 Wisp\\

\noindent \textbf{Duration:} 1 scene\\

\noindent The Princess forms a Wisp into an illusory object. The illusion affects only sight, and material objects pass through it without stopping. The Size of the image is limited to 2 + Inner Light or any applied Invocation, but within that limit she can show any image she is capable of imagining.

\begin{myindentpar}{10pt}

\emph{Dramatic Failure:} The Charm fails with a burst of light and a loud crack, alerting anyone nearby who could see or hear to the Princess's presence.

\emph{Failure:} No image appears.

\emph{Success:} The Princess creates the image she desires, somewhere within her sensory range. The image remains until the Charm ends. The Princess may manipulate the image for no cost, for example she could make an illusory person walk around but could not turn it into an illusory car. This requires her full attention, if a roll is required the dice pool is Wits + Expression though the Storyteller may nominate a different dice pool if the situation warrants it. E.G. Wits + Subterfuge to make an illusory person deceive someone. Intelligence + Academics to create an illusory chalkboard with accurate scientific equations. 

\emph{Exceptional Success:} The Princess manipulates her illusion deftly. Her mundane actions to control it gain a +2 bonus.

\textbf{Upgrade: Dancing}

The Princess may create an image that moves without her conscious attention. When her concentration lapses, the image continues to act on the instructions she last gave it.

\textbf{Upgrade: Glowing}

Requires Light the Way

The Princess can make her illusions give off light, just as her Regalia does with Light the Way. As long as she is concentrating on an image, she may vary the light coming from it reflexively, independently of her personal light. If she has upgraded Light the Way, she may apply all that Charm’s upgrades to the image’s light as if it were her own.

\textbf{Upgrade: Consecrated}

The Princess may sustain a illusion upon the blessing of a Consecrated area. Doing to keeps the illusion intact until the Consecrated Condition ends, but the illusion can never manifest beyond the area covered by the Condition. 

\textbf{Upgrade: Singing}

The Princess may create an illusory sound, or add sound to an illusory image. A sound without an image seems to come from any point the Princess can see, which she can move while concentrating.

\textbf{Upgrade: Solid}

The image becomes substantial. As long as it lasts it has a degree of solidity, as well as a firmness and texture appropriate to what it depicts. However, just because it projects an image of solidity does not mean it is solid; the image shatters if force is applied to it (by kicking an illusory wall, for example) or with it (an illusionary wrench won't tighten or loosen a bolt).

\textbf{Upgrade: Thinking (Acqua \textbullet\textbullet)}

Requires Dancing

When the Princess releases her concentration, in addition to providing the illusion with further instructions she may grant it a measure of adaptability. For each dot of Acqua she may instruct the image to recognize one condition and react in response. The illusion has mundane senses within the human range. If a roll is required the illusion has a dice pool equal to the Princess' Acqua, for an illusory watchman who sounds the alarm when someone approaches and could roll Acqua to contest stealth.

\textbf{Upgrade: Multiplied (Aria \textbullet\textbullet)}

Requires Dancing

The Princess may create multiple distinct images with a single Charm activation. She can make one image for each dot of her Aria; for each image beyond the first, the maximum Size of all the images is reduced by 1. The Princess can try to control more than one image at once but takes a -2 penalty to her dice pool for each image beyond the first. If she applies Singing, each image may have its own sound and images can be replaced with sounds, but the number of distinct illusions remains limited to her Aria.

\textbf{Upgrade: Scented (Legno \textbullet)}

The Princess may create an illusory scent, or add scent to an illusory image or sound. If the image is Solid, she can give it a taste as well. A scent without an image seems to come from a point the Princess can see, which she can move while concentrating.


\end{myindentpar}

\subsubsection*{Face of Lover's Alarm (Appear \textbullet\textbullet, Specchio \textbullet\textbullet)}

\noindent \textbf{Action:} Reflexive

\noindent \textbf{Dice Pool:} Manipulation + Empathy vs Wits + Composure

\noindent \textbf{Cost:} 1 Wisp\\

\noindent One would not threaten a loved one, would they? And as the Queen of Mirrors teaches, the True Heir will be beloved by all. This Charm is activated in response to an attack before the target rolls and the Princess must be able to apply her Defense to the attack.

\begin{myindentpar}{10pt}\emph{Dramatic Failure:} The target sees the Noble as a perverted mockery of one they love, and gain +3 dice on their attack roll. They cannot be affected by this Charm for the reminder of the scene.

    \emph{Failure:} The target sees through the illusion and cannot be affected by this Charm for the reminder of the scene.

    \emph{Success:} The Noble's face shifts into someone the target loves and they are compelled to react appropriately, pulling their blow aside before the Princess can be hurt. The wasted attack still consumes their action for the turn. They also take the Blessing Condition and gain +3 dice to resist this Charm for the remainder of the scene. 

    \emph{Exceptional Success:}  As above, and the target takes some small Tilt. They might unbalance themselves and take Knocked Down.  

\textbf{Upgrade: Inviolate (Specchio \textbullet\textbullet\textbullet)}

\noindent \textbf{Cost:} +1 Willpower

Just as the Princess shifts into someone the target loves, someone else in reach of the target's weapon takes the Princess' form. The Princess nominates another person within range of the attacker's weapon, if she wins the Contested roll the target must attack them instead. 

\textbf{Upgrade: Doomed (Specchio \textbullet\textbullet\textbullet\textbullet)}

\noindent \textbf{Cost:} +1 Willpower

Does the mirror reflect the truth, or does it show things how they should be? If the individual whom the Noble mimics is present, and within the attack range of the form of weapon used to attack her, the attack is transferred to them. If the Princess wins the contested roll the target must make an attack roll against the her, subtracting her Defense as normal, but the damage goes to their loved one. In the case of a melee weapon, the target might swing wide and hit their friend; for a ranged weapon, it can be more improbable, as ricochets, arcing electrical conduits, and exploding fire extinguishers all conspire to hurt their ally. If the loved one has armor or similar protections they apply as normal. \end{myindentpar}

\subsubsection*{Speaking Eyes (Appear \textbullet\textbullet, Acqua \textbullet)}
\noindent \textbf{Action:} Instant

\noindent \textbf{Dice Pool:} Wits + Persuasion

\noindent \textbf{Cost:} 1 Wisp

\noindent \textbf{Duration:} 1 scene\\

\noindent The movements of the Princess' body are saturated by meaning. She uses this Charm on a single individual who can see her. 

\begin{myindentpar}{10pt}

\emph{Dramatic Failure:} The Noble destroys her natural ability to communicate. All her Expression rolls take a -2 penalty for the rest of the scene.

\emph{Failure:} The Noble’s gestures convey no special meaning.

\emph{Success:} The Noble can convey the most complicated thoughts to her target with a turn of the head, a twist of the hand, or a facial expression, as quickly as with normal speech; the target will understand her intent even if he shares no languages with her, as long as they can see her. Nobody else looking at her will understand her meaning; indeed, unless onlookers examine her carefully (succeeding on an Intelligence + Investigation roll) they won’t realize that she is “saying” anything at all. The Charm does not help the Noble understand the target.

\emph{Exceptional Success:}  The Noble may combine her gestures with her speech. Her Expression rolls to communicate with the target gain a +2 bonus until the Charm ends.

\textbf{Upgrade: Collective (Acqua \textbullet\textbullet)}

Modified by Commonalty

The Princess can convey her thoughts to members of an organization that can see her, using the Commonalty modifier. Those not in the organization will not understand her gestures.\end{myindentpar}  

\subsubsection*{Words on Water Writ (Appear \textbullet\textbullet, Acqua \textbullet)}

\noindent \textbf{Action:} Instant

\noindent \textbf{Dice Pool:} Intelligence + Expression - Commonality

\noindent \textbf{Cost:} 1 Wisp

\noindent \textbf{Duration:} Lasting\\

\noindent The Princess encodes a message into a seemingly innocuous display. First she must create a physical display to contain her message, and it could be anything. She could arrange her collection of stuffed toys into a symbolic order, or write a letter with a second coded message. Then she names the recipient or recipients. She may choose one person, a group (applying a Commonality modifier), a Court (-5) or the Nobility as a whole (-5). 

\begin{myindentpar}{10pt}

\emph{Dramatic Failure:} The Princess message is not hidden as well as she thinks. At some point in the future the Storyteller should introduce a problem caused by it falling into the wrong hands. 

\emph{Failure:} The charm fails to activate.

\emph{Success:} The Princess creates a message which contains roughly as much information as one minute's worth of speech per success. Regular speech, not the Royal Tongue. When any of the intended recipients examines the object containing they immediately know it is significant. No roll is necessary to decode the message but some time and attention might be. If the object is destroyed or significantly modified or damaged the hidden message is lost. If the message falls into the hands of anyone but the intended recipient they will not be able to decrypt it without magic and a Clash of Wills.  

\emph{Exceptional Success:}  The Princess may optionally add a duration to the magic, just encase it falls into the hands of a magical opponent. \end{myindentpar}


\subsubsection*{Summon Backup Dancers (Appear \textbullet\textbullet\textbullet)}
\noindent \textbf{Action:} Instant

\noindent \textbf{Dice Pool:} Dexterity + Expression

\noindent Cost 3 Wisps

\noindent \textbf{Duration:} 1 scene \\

\noindent Sometimes a Princess just needs to wow a crowd. With this Charm a Princess can summon her own support team with a maximum of Inner Light members.

Before casting Summon Backup Dancers the Princess must declare what kind of action they're going to support. The Princess could declare singing in which case they'd be able to perform dance routines or backup vocals. The Princess could declare painting and summon a model to paint. The backup dancers primary task must involve giving some sort of performance, she cannot summon a team of stage hands and set them wiring up her stage. 

Backup dancers are solid but only when performing the task they were created, an actor would be able to interact with props, and wear costumes, but would not be able to fetch a cup of coffee. A cunning Princess might still be able to use them for other roles, she might ask her trope to dance around waving their arms in front of her enemies face while she hides, or to attack foes who don't realize the backup dancers cannot hurt them. (Assume a 1 die pool on any social actions, mental and physical tasks fail automatically). If attacked a single point of Damage, even Bashing, causes a Backup dancer to vanish into a puff of ephemera appropriate to the Princess.

A Princess can only have one instance of Summon Backup Dancers active at a time.

\begin{myindentpar}{10pt}     \emph{Dramatic Failure:} You summon Inner Light performers but have no control over them. They immediately run riot, while not actually harmful they can be very very annoying and are certainly inappropriate to the Princesses wishes. Think clowns, not monsters.

    \emph{Failure:} No Dancers appear.

    \emph{Successes:} The Noble creates several assistants. The crew acts as equipment for the chosen task, giving an equipment bonus equal to the Successes rolled (maximum of 5) to the Noble, or to anyone else she tells them to help. However, they carry the Volatile Condition [CofD 102]. If the Princess performs badly her backup will follow her lead, potentially following her off a metaphorical cliff. 

    \emph{Exceptional Success:} he crew doesn’t flake out when their principle blunders. They don’t carry the Volatile Condition \end{myindentpar}

\subsubsection*{Goccia Astrale (\textbullet\textbullet\textbullet)}
\noindent \textbf{Action:} Instant

\noindent \textbf{Dice Pool:} Intelligence + Crafts 

\noindent \textbf{Cost:} 2+ Wisps

\noindent \textbf{Duration:} Variable\\ 

\noindent \textquotedblleft I can't be in two places at once!\textquotedblright\, is an all-too-common complaint from just about anyone; the Hopeful especially find it necessary to be two people at once, when mundane obligations interfere with fighting a supernatural incursion. With this Charm a Princess can partly resolve such dilemmas, by making a facsimile of her mundane self that moves and speaks, and can pass as herself if not examined too closely.

To use the Charm, the Princess must have a mass of material that she can shape barehanded, or else a quantity of fluid, of roughly Size 2. Infusing this material with Wisps, she transforms it into an Astral Droplet: an exact physical copy of herself just before her last transformation, wearing identical clothing. The Droplet has the same mundane Attributes and Skills as its creator, and access to all her memories as of the time of its creation. However, it lacks energy and imagination; it does not get 10-again on any dice pool, or any other quality that allows rerolling dice, cannot spend Willpower, and has neither Virtue nor Vice. The Droplet is also fragile - a single point of damage of any type is enough to dissolve it into the stuff it was made of and glimmers of light. Finally, the Droplet has none of the Princess' Light-derived powers, not even Practical Magic.

The Droplet's basic duration is 30 minutes for each activation success. The Princess may choose to extend the duration by spending more Wisps than the minimum of 2 - each extra Wisp doubles the time before the Droplet disintegrates, cumulatively. She cannot, however, spend more Wisps for this purpose than half her Inner Light, rounded up. She may choose the total number of Wisps spent after the activation roll, though she must spend at least 2. She may end the Charm prematurely by touching the Droplet and concentrating for a turn. A Princess cannot create a second Droplet before the first Droplet dissolves - the Charm just fails if she tries.

By default the Droplet is an automaton which follows any instruction given to it which does not pose a threat of physical harm, and otherwise stays wherever it is put; it takes no action on its own. If the material used to make it resonates with an Invocation the Princess may apply that Invocation to the Charm (Aria uses thick mist or smoke, or a strong breeze, Lacrima can use total darkness or dark shadow.) Droplets so formed have some initiative and the rudiments of a personality; in the absence of clear instructions they follow the tenets of the applied Invocation, to the extent allowed by the knowledge and abilities the Princess possesses.

\begin{myindentpar}{10pt} \textbf{Upgrade: Integrated}

It's often useful to know what your double has been doing during your absence. If the Princess dissolves her Droplet prematurely by touching it and concentrating, she can recall to memory a description of the actions the Droplet took and the events it saw and heard. A roll to recall any detail of the account is at -2 in addition to any other penalties.

\textbf{Upgrade: Lifelike (Specchio \textbullet\textbullet\textbullet)}

A Lightbringer can use that indescribable quality that makes her reflection \textquotedblleft her reflection\textquotedblright instead of a simple image and use it to place mind and will into the Droplets they make. To apply Lifelike the Princess needs a reflecting surface large enough to hold an image of her upper body; she takes the hands of her image and pulls it out of the surface. The resulting Droplet is flipped left-to-right, and with the opposite dominant hand. The real difference though, is that the Droplet gains a nearly-human initiative and will; it has the Belief dots, Virtues and Vice of its creator, gets 10-again on its dice pools, can benefit from other qualities that permit rerolls, and may spend and regain Willpower (though it begins with none.) The only supernatural ability the Droplet possesses is the ability to Transform and benefit from Transformed skills and attributes, created by Specchio it possesses the Queen of Mirror's Practical Magic. The first time the Princess creates this 
Droplet, its personality exactly matches her own (including any forms of madness) and aside from the physical inversion they can't be told apart.

\textbf{Drawback:} By invoking a Lifelike a Princess is summoning an independent being. The Droplet's personality does not vanish when its body dissolves; it returns to being her reflection, until she makes another Lifelike Droplet. At that moment it returns and inhabits the new-made body. Since the droplet is created from a reflection it sees nothing wrong with going back behind the mirror (and finds acting independently a bit weird and uncomfortable, but nobody said being a Princess or her reflection was easy). However every time the Princess invokes Lifelike there is a chance that Specchio's madness will infect her Droplet. Each time it inhabits a new body the Droplet rolls Resolve + Composure, with a -2 penalty if the Princess is suffering from Specchio induced madness. If she fails, the Droplet's falls to Specchio's arrogance; she is convinced of her excellence and her right to rule, and resents the Princess taking her rightful limelight. The Droplet will do what it can to avoid returning to behind the 
mirror. It can lengthen the Charm's duration by spending its own Willpower; it may spend 1 point each day for this, which adds 24 hours of continued embodiment. (The Droplet cannot do this if the Resolve + Composure check succeeds.) 

Just what the Droplet does is up to the Storyteller, but it should avoid Compromising the Princess' own Beliefs.

If the Princess applies Integrated and Lifelike, all her memories since the Droplet's first creation become available to it on the same terms as its are available to her after she dissolves it. Other than that, neither personality can remember what the other has done.\end{myindentpar}

\subsubsection*{Wallflower (Appear \textbullet\textbullet\textbullet)}
\noindent \textbf{Action:} Instant

\noindent \textbf{Dice Pool:} Manipulation + Stealth

\noindent \textbf{Cost:} 1 Wisp

\noindent \textbf{Duration:} 1 scene \\

\noindent The Princess shies out of view entirely. She becomes both physically invisible and mentally invisible as people find themselves incapable of focusing on evidence of her activities. If she opens a door and walks through nobody will notice the door move, though they may unconsciously move to avoid being hit by the door. 

Any use of Charms ends this power, as does Swords or Tears' Practical Magic. Any act of violence, or any act that could be described as \textquotedblleft attention grabbing\textquotedblright\, such as yelling or jumping up and down while waving your arms ends the effect. 

This protection extends to anyone viewing a recording or photo of the Princess that is both made and viewed during the current scene. However once the Charm ends so does her psychic invisibility, meaning that evidence of her activity will no longer be ignored. If she is caught on camera opening a door then once the Charm ends anyone reviewing the footage will see a door mysteriously open by itself. 

\begin{myindentpar}{10pt}

\emph{Dramatic Failure:} The Noble makes herself memorable without realizing it. She takes the Inverse Curse Condition, all Perception rolls to notice her take a +2 bonus for the rest of the scene.

\emph{Failure:} Nothing happens. 

\emph{Success:} The Princess stops emitting any light in the visible spectrum and clouds the minds of any who might notice her regardless. She is essentially impossible to notice with mundane senses. Supernatural senses allow a Clash of Wills roll to penetrate the Princess' invisibility. A mortal who suspects the Princess is present my track her by looking for footsteps, opening doors, or other evidence of her passing. This requires an Extended Wits + Investigation roll with a target of the Princess' Composure + Stealth + Inner Light. Even with useful equipment like an infra-red camera it will take a while for the tracker to force his mind to except that heat patch is the Princess and not a glitch or passer by. 

\emph{Exceptional Success:} The Princess can shout and wave her arms as much as she wants. Unless she intends for someone to notice her the Charm will not end. Violence, Charms and Swords or Tear's practical magic will still end the Charm. 

\textbf{Upgrade: Collective}

The Princess can hide an entire group by applying the Commonality modifier. 

\textbf{Upgrade: Veiled}

Using another Charm does not automatically end this Charm, though the results of the Charm are not masked, neither are any flashes or sparkles created as a side effect.

\textbf{Upgrade: Full Spectrum (Acqua \textbullet\textbullet)}

The Princess' invisibility extends into the infra-red and ultra-violet spectrum. In addition the mental effects of the Charm extend to affect mundane technology, making her all but invisible to all technological means. A door that is set to issue an alarm when it is opened will remain silent unless the Princess leaves it open when the charm ends or someone else jostles it. 
\end{myindentpar}

\subsubsection*{A Flight of Furious Fancies (Appear \textbullet\textbullet\textbullet, Aria \textbullet\textbullet\textbullet)}

\noindent \textbf{Action:} Instant

\noindent \textbf{Dice Pool:} Wits + Subterfuge - Composure

\noindent \textbf{Cost:} 1 Wisp

\noindent \textbf{Duration:} 1 Turn\\

\noindent The princess conjures an host of illusory attacks on any target she can see clearly. Their attacks can do no damage, but the target may waste their attention defending against non-existent foes.

\begin{myindentpar}{10pt}     \emph{Dramatic Failure:} The swirling of the host leaves the Princess dizzy and confused. She immediately takes the Stunned Tilt [CofD 286].

    \emph{Failure:} No attackers appear.

    \emph{Successes:}  The target is thoroughly distracted by the host’s attacks. Each success reduces the target's Defense by two for the rest of the turn.

    \emph{Exceptional Success:} In addition the target takes the Stunned Tilt. \end{myindentpar}


\subsubsection*{Doppelgänger Defense (Appear \textbullet\textbullet\textbullet, Aria \textbullet\textbullet\textbullet)}
\noindent \textbf{Action:} Reflexive

\noindent \textbf{Dice Pool:} Wits + Subterfuge vs Wits + Composure

\noindent \textbf{Cost:} 2 Wisps, 1 Willpower

\noindent \textbf{Duration:} 1 Turn\\

\noindent At the perfect moment the princess creates an illusory double of her ally, or herself, just as her enemy prepares to strike. If the Princess wins the contested roll her opponent targets the illusion, if she fails her enemy sees through the ruse or gets lucky.

\begin{myindentpar}{10pt}     \emph{Dramatic Failure:} The Princess is unable to tell herself and the illusion apart. She immediately takes the Stunned Tilt [CofD 286].

    \emph{Failure:} No illusions appear.

    \emph{Successes:}  The target is fooled by the illusion. The effects of this Charm depend on what the target was planning to do. A sword stroke is entirely wasted but throwing a grenade at the illusion might still catch the real person in the blast. If the target is using a gun capable of Autofire she could easily shoot at both to be on the safe side.

    \emph{Exceptional Success:} In addition the illusion persists for a moment longer. The Princess may invoke this Charm a second time before the end of the turn without paying Wisps or Willpower. \end{myindentpar}




\subsubsection*{Know My Pain (Appear \textbullet\textbullet\textbullet, Tempesta \textbullet\textbullet)}
\noindent \textbf{Action:} Instant

\noindent \textbf{Dice Pool:} Resolve + Intimidation vs Composure + Supernatural Advantage

\noindent \textbf{Cost:} 2 Wisps

\noindent \textbf{Duration:} 1 scene\\ 

\noindent There are those who would tell you that pain is bad, and to be avoided at all costs. They are weak. Pain is a tool, just like any other. It is not only useful, but even necessary, to use it against those who would try to stop you.

\begin{myindentpar}{10pt}    \emph{Dramatic Failure:} The Princess gives pain only to herself. For the rest of the scene she suffers a -1 penalty to all her actions.

    \emph{Failure:} The target ignores the Princess' pain.

    \emph{Success:} The target suffers phantom bashing damage, equal to the number of boxes with damage of any type on the Princess' Health track or his unmarked Health boxes, whichever is less. The target suffers wound penalties from this damage if it goes into the last three boxes, and must check their Stamina to remain conscious if it fills the last box, as normal. However, the damage is not real; it vanishes when the Charm ends as if it were healed, and if the target takes real bashing damage when the track is full, that damage replaces the phantom damage instead of wrapping into lethal damage. Phantom damage also has no effect if the target has no empty Health Boxes to fill. The phantom damage can also be healed magically, or mitigated as it's inflicted by targets who have that ability.

    \emph{Exceptional Success:} The phantom damage from the Charm is lethal. The target does not bleed out if the phantom damage incapacitates him, and real lethal damage replaces it instead of wrapping into aggravated damage. 

\textbf{Upgrade: Shared (Tempesta \textbullet\textbullet\textbullet)}

\noindent \textbf{Cost:} +1 Wisp

\noindent Modified by Commonalty

The Princess may use Know My Pain on a social group, applying the Commonalty modifier; every member of the group takes phantom damage. The member with the highest Composure resists for the whole group.

\end{myindentpar}

\subsubsection*{Imaginary Friend (Appear \textbullet\textbullet\textbullet\textbullet, Aria \textbullet\textbullet\textbullet)}

\noindent \textbf{Action:} Instant, Manipulation + Subterfuge vs. Wits + Subterfuge + Supernatural Advantage

\noindent \textbf{Cost:} 2 Wisps

\noindent \textbf{Duration:} 1 scene\\

\noindent The Noble creates the illusion of a man or woman, especially crafted for a single person.

\begin{myindentpar}{10pt}    \emph{Dramatic Failure:} The target becomes deeply suspicious of the Noble; for the rest of the scene all her Social rolls aimed at him take a -2 penalty.

    \emph{Failure:}  The Noble fails to create the illusion.

    \emph{Success:} The Noble chooses a role for the illusion; the target perceives a person who behaves appropriately for that role in the present situation, according to the target’s beliefs. However, if she doesn’t beat the target’s successes, he realizes that this person is an illusion within moments. If she does get more successes, the target will treat the imaginary person as fully real, and respond to it as to a person of the role the Noble chose. No one else - including the Noble - will see or hear the illusion; she has no control over its behaviour.

    \emph{Exceptional Success:} The Noble can see the illusion. 	\end{myindentpar}

\subsubsection*{Garden of Bright Images (Appear \textbullet\textbullet\textbullet\textbullet)}
\noindent Requires: Phantom

\noindent \textbf{Action:} Extended, Presence + Expression (1 minute/roll, threshold = Sanctuary + duration in hour). 

\noindent \textbf{Cost:} (Sanctuary rating) Wisps

\noindent \textbf{Duration:} 1 or more hours \\

\noindent The Princess decorates her surroundings with a beautiful glamour. She chooses an area to affect and a time for the glamour to last when she starts working on the Charm.

\begin{myindentpar}{10pt}    \emph{Dramatic Failure:} The Princess deceives only herself; she sees the area transformed according to her intentions, but everyone else sees it as it really is

\emph{Failure:} The Princess makes no progress towards the Threshold.

\emph{Success:} When the Princess reaches the threshold, the appearance of every nonliving thing in the targeted area is transformed to match a theme the Princess chooses. (For instance, \textquotedblleft formal ballroom", \textquotedblleft open bazaar", \textquotedblleft African jungle" and \textquotedblleft undersea grotto" are all valid themes.) Permanent fixtures are always changed; she can choose whether to change people’s possessions, but cannot affect some possessions and not others. As long as the Charm lasts, the changed objects look, sound, feel and even smell like what they appear to be; observers may make Perception rolls to spot something wrong with the illusion, but they take a penalty equal to the Princess’ successes.

\emph{Exceptional Success:} If the Princess scores an Exceptional Success on her extended roll then mundane senses cannot penetrate the illusion. 

\textbf{Upgrade: Palatine} 

The Princess may take advantage of a Consecrated Condition. When she does so, the illusion covers all the Consecrated ground and lasts until the Condition expires or she dismisses the Charm; treat the duration as 0 when calculating the threshold\end{myindentpar}